\documentclass[12pt,a4paper,twoside,openright]{report}
\let\openright=\cleardoublepage



%%% Choose a language %%%

\newif\ifEN
\ENtrue   % uncomment this for english
%\ENfalse   % uncomment this for czech

%%% Configuration of the title page %%%

\def\ThesisTitleStyle{mff} % MFF style
%\def\ThesisTitleStyle{cuni} % uncomment for old-style with cuni.cz logo
%\def\ThesisTitleStyle{natur} % uncomment for nature faculty logo

\def\UKFaculty{Faculty of Mathematics and Physics}
%\def\UKFaculty{Faculty of Science}

\def\UKName{Charles University in Prague} % this is not used in the "mff" style

% Thesis type names, as used in several places in the title
%\def\ThesisTypeTitle{\ifEN BACHELOR THESIS \else BAKALÁŘSKÁ PRÁCE \fi}
\def\ThesisTypeTitle{\ifEN MASTER THESIS \else DIPLOMOVÁ PRÁCE \fi}
%\def\ThesisTypeTitle{\ifEN RIGOROUS THESIS \else RIGORÓZNÍ PRÁCE \fi}
%\def\ThesisTypeTitle{\ifEN DOCTORAL THESIS \else DISERTAČNÍ PRÁCE \fi}
%\def\ThesisGenitive{\ifEN bachelor \else bakal{\'a}\v{r}sk{\'e} \fi}
\def\ThesisGenitive{\ifEN master \else diplomov{\'e} \fi}
%\def\ThesisGenitive{\ifEN rigorous \else rigor{\'o}zn{\'i} \fi}
%\def\ThesisGenitive{\ifEN doctoral \else diserta\v{c}n{\'i} \fi}
%\def\ThesisAccusative{\ifEN bachelor \else bakal{\'a}\v{r}skou \fi}
\def\ThesisAccusative{\ifEN master \else diplomovou \fi}
%\def\ThesisAccusative{\ifEN rigorous \else rigor{\'o}zn{\'i} \fi}
%\def\ThesisAccusative{\ifEN doctoral \else diserta\v{c}n{\'i} \fi}



%%% Fill in your details %%%

\def\ThesisTitle{Inference-driven resource managemenent and polymorphism in systems programming}
\def\ThesisAuthor{Ji\v{r}{\'i} Klepl}
\def\YearSubmitted{2022}

% department assigned to the thesis
\def\Department{Department of Software Engineering}
% Is it a department (katedra), or an institute (ústav)?
\def\DeptType{Department}

\def\Supervisor{RNDr. Miroslav Kratochv{\'i}l, Ph.D.}
\def\SupervisorsDepartment{Department of Software Engineering}

% Study programme and specialization
\def\StudyProgramme{Computer Science - Software Systems}
\def\StudyBranch{ISWSP}

\def\Dedication{%
I would like to thank my supervisor and mentor Mirek Kratochv{\'i}l for his support and guidance throughout my work on this thesis. His input was valuable in all aspects of the work, including the intricacies of academic writing and typesetting. I~would also like to thank my family and friends for supporting me and encouraging me during my studies.
}

\def\AbstractEN{%
Systems programming languages facilitate the implementation of software that runs in restricted environments close to the hardware, such as operating systems, drivers, and real-time and embedded systems.
Implementation of desirable features of such languages, such as generic programming and type system capabilities that prevent programmer errors, are complicated by strict constraints on the run-time properties of the program.
This thesis explores a novel combination of C-\phantom{}- language with advanced type system features that allow type checking of highly polymorphic generic code and demonstrates type-driven resource management in this language.
As the main result, the thesis provides a proof-of-concept in a prototype compiler of an extended version of C-\phantom{}- to LLVM and describes a type system based on deferred constraint solving that is capable of type inference in the presence of multi-parameter typeclasses and C-\phantom{}- subtypes.
We demonstrate the functionality of the type system and the compiler on selected program examples and report several identified design challenges that may be addressed to make the system more practical.
}

\def\AbstractCS{%
Syst{\'e}mov{\'e} programovac{\'i} jazyky usnad\v{n}uj{\'i} implementaci softwaru, kter{\'y} b\v{e}\v{z}{\'i} v omezen{\'y}ch prost\v{r}ed{\'i}ch bl{\'i}zko hardwaru, jako jsou opera\v{c}n{\'i} syst{\'e}my, drivery, syst{\'e}my re{\'a}ln{\'e}ho \v{c}asu a vestav\v{e}n{\'e} syst{\'e}my.
Implementace prosp\v{e}\v{s}n{\'y}ch prvk\r{u} pro tyto jazyky, jako jsou generick{\'e} programov{\'a}n{\'i} a typov{\'y} syst{\'e}m zaji\v{s}\v{t}uj{\'i}c{\'i} bezchybnost k{\'o}du, je zt\v{e}\v{z}ov{\'a}na úzk{\'y}mi po\v{z}adavky na b\v{e}hov{\'y}m{\'e} vlastnostni programu.
Hlavn{\'i} n{\'a}pln{\'i} t{\'e}to pr{\'a}ce je prozkoumat novou kombinaci jazyka C-\phantom{}- s prvky typov{\'e}ho syst{\'e}mu umo\v{z}\v{n}uj{\'i}c{\'i} typovou kontrolu vysoce polymorfn{\'i}ho k{\'o}du a demonstruje mo\v{z}nosti typov\v{e} veden{\'e} spr{\'a}vy zdroj\r{u} v tomto jazyku.
Jako hlavn{\'i} v{\'y}stup pr{\'a}ce pod{\'a}v{\'a} d\r{u}kaz konceptu zprost\v{r}edkovan{\'y} prototypem p\v{r}eklada\v{c}e ro\v{z}\v{s}{\'i}\v{r}en{\'e}ho jazyka C-\phantom{}- do LLVM a popisuje typov{\'y} syst{\'e}m zalo\v{z}en{\'y} na odkl{\'a}daj{\'i}c{\'i}m \v{r}e\v{s}en{\'i} omezuj{\'i}c{\'i}ch podm{\'i}nek, kter{\'y} je schopen typov{\'e} inference za p\v{r}{\'i}tomnosti v{\'i}ce-parametrov{\'y}ch typov{\'y}ch t\v{r}{\'i}d a subtypy C-\phantom{}-.
Funkcionalitu typov{\'e}ho syst{\'e}mu a p\v{r}eklada\v{c}e demonstrujeme na vybran{\'y}ch uk{\'a}zk{\'a}ch program\r{u} a popisujeme n\v{e}kolik rozpoznan{\'y}ch komplikac{\'i} n{\'a}vrhu, kter{\'e} lze adresovat pro prakti\v{c}t\v{e}j\v{s}{\'i} implementaci.
}

% 3 to 5 keywords (recommended), each enclosed in curly braces.
% Keywords are useful for indexing and searching for the theses by topic.
\def\Keywords{%
{programming languages}, {type systems}, {compilers}, {systems programming}, {resource management}
}

% If your abstracts are long and do not fit in the infopage, you can make the
% fonts a bit smaller by this setting. (Also, you should try to compress your abstract more.)
% Alternatively, consider increasing the size of the page by uncommenting the
% geometry modification in thesis.tex.
\def\InfoPageFont{}
%\def\InfoPageFont{\small}  %uncomment to decrease font size

\ifEN\relax\else
% If you are writing a czech thesis, you additionally need to fill in the
% english translation of the metadata here!
\def\ThesisTitleEN{Inference-driven resource managemenent and polymorphism in systems programming}
\def\DepartmentEN{Department of Software Engineering}
\def\DeptTypeEN{Department}
\def\SupervisorsDepartmentEN{Department of Software Engineering}
\def\StudyProgrammeEN{Computer Science - Software Systems}
\def\StudyBranchEN{ISWSP}
\def\KeywordsEN{%
{programming languages}, {type systems}, {compilers}, {systems programming}, {resource management}
}
\fi


\usepackage[a-2u]{pdfx}

\ifEN\else\usepackage[czech,shorthands=off]{babel}\fi
\usepackage[utf8]{inputenc}
\usepackage[T1]{fontenc}

% See https://en.wikipedia.org/wiki/Canons_of_page_construction before
% modifying the size of printable area. LaTeX defaults are great.
% If you feel it would help anything, you can enlarge the printable area a bit:
%\usepackage[textwidth=390pt,textheight=630pt]{geometry}
% The official recommendation expands the area quite a bit (looks pretty harsh):
%\usepackage[textwidth=145mm,textheight=247mm]{geometry}

%%% FONTS %%%
\usepackage{lmodern} % TeX "original" (this sets up the latin mono)

% Optionally choose an override for the main font for typesetting
\usepackage[mono=false]{libertinus} % popular for comp-sci (ACM uses this)
%\usepackage{tgschola} % Schoolbook-like (gives a bit of historic feel)
%\usepackage[scale=0.96]{tgpagella} % Palladio-like (popular in formal logic).

% Optionally choose a custom sans-serif fonts (e.g. for figures and tables).
% Default sans-serif font is usually Latin Modern Sans. Some font packages
% (e.g. libertinus) replace that with a better matching sans-serif font.
%\usepackage{tgheros} % recommended and very readable (Helvetica-like)
%\usepackage{FiraSans} % looks great
% DO NOT typeset the main text in sans-serif font!
% The serifs make the text easily readable on the paper.

% IMPORTANT FONT NOTE: Some fonts require additional PDF/A conversion using
% the pdfa.sh script. These currently include only 'tgpagella'; but various
% other fonts from the texlive distribution need that too (mainly the Droid
% font family).


% some useful packages
\usepackage{microtype}
\usepackage{amsmath,amsfonts,amsthm,bm,harpoon}
\usepackage{graphicx}
\usepackage{xcolor}
\usepackage{booktabs}
\usepackage{caption}
\usepackage{floatrow}

% load bibliography tools
\usepackage[backend=bibtex,natbib,style=numeric,sorting=none]{biblatex}
% alternative with alphanumeric citations (more informative than numbers):
%\usepackage[backend=bibtex,natbib,style=alphabetic]{biblatex}
%
% alternatives that conform to iso690
% (iso690 is not formally required on MFF, but may help elsewhere):
%\usepackage[backend=bibtex,natbib,style=iso-numeric,sorting=none]{biblatex}
%\usepackage[backend=bibtex,natbib,style=iso-alphabetic]{biblatex}
%
% additional option choices:
%  - add `giveninits=true` to typeset "E. A. Poe" instead of full Edgar Allan
%  - `terseinits=true` additionaly shortens it to nature-like "Poe EA"
%  - add `maxnames=10` to limit (or loosen) the maximum number of authors in
%    bibliography entry before shortening to `et al.` (useful when referring to
%    book collections that may have hundreds of authors)
%  - for additional flexibility (e.g. multiple reference sections, etc.),
%    remove `backend=bibtex` and compile with `biber` instead of `bibtex` (see
%    Makefile)
%  - `sorting=none` causes the bibliography list to be ordered by the order of
%    citation as they appear in the text, which is usually the desired behavior
%    with numeric citations. Additionally you can use a style like
%    `numeric-comp` that compresses the long lists of citations such as
%    [1,2,3,4,5,6,7,8] to simpler [1--8]. This is especially useful if you plan
%    to add tremendous amounts of citations, as usual in life sciences and
%    bioinformatics.
%  - if you don't like the "In:" appearing in the bibliography, use the
%    extended style (`ext-numeric` or `ext-alphabetic`), and add option
%    `articlein=false`.
%
% possibly reverse the names of the authors with the default styles:
%\DeclareNameAlias{default}{family-given}

% load the file with bibliography entries
\addbibresource{refs}

% remove this if you won't use fancy verbatim environments
\usepackage{fancyvrb}

% remove this if you won't typeset TikZ graphics
\usepackage{tikz}
\usetikzlibrary{positioning} %add libraries as needed (shapes, decorations, ...)

% remove this if you won't typeset any pseudocode
\usepackage{algpseudocode}
\usepackage{algorithm}
\usepackage{array}

% remove this if you won't list any source code
\usepackage{listings}

% languages
\usepackage{formal-grammar}

\hypersetup{unicode}
\hypersetup{breaklinks=true}

\usepackage[noabbrev]{cleveref}


% various forms of TODOs (you should remove this before submitting)
\usepackage[textsize=tiny, backgroundcolor=yellow!25, linecolor=black!25]{todonotes}
\newcommand{\xxx}[1]{\textcolor{red!}{#1}}

 % remove this before compiling the final version


% use this for typesetting a chapter without a number, e.g. intro and outro
\def\chapwithtoc#1{
\chapter*{#1}
\addcontentsline{toc}{chapter}{#1}
}

% If there is a line/figure overflowing into page margin, this will make the
% problem evident by drawing a thick black line at the overflowing spot. You
% should not disable this.
\overfullrule=3mm

% The maximum stretching of a space. Increasing this makes the text a bit more
% sloppy, but may prevent the overflows by moving words to next line.
\emergencystretch=1em

\ifEN
\theoremstyle{plain}
\newtheorem{thm}{Theorem}
\newtheorem{lemma}[thm]{Lemma}
\newtheorem{claim}[thm]{Claim}
\newtheorem{defn}{Definition}
\newtheorem{ex}{Example}
\theoremstyle{remark}
\newtheorem*{cor}{Corollary}
\else
\theoremstyle{plain}
\newtheorem{thm}{Věta}
\newtheorem{lemma}{Lemma}
\newtheorem{claim}{Tvrzení}
\newtheorem{defn}{Definice}
\newtheorem{ex}{Ukázka}
\theoremstyle{remark}
\newtheorem*{cor}{Důsledek}
\fi

\newenvironment{myproof}{
  \par\medskip\noindent
  \textit{\ifEN Proof \else Důkaz \fi}.
}{
\newline
\rightline{$\qedsymbol$}
}

% real/natural numbers
\newcommand{\R}{\mathbb{R}}
\newcommand{\N}{\mathbb{N}}

% asymptotic complexity
\newcommand{\asy}[1]{\mathcal{O}(#1)}

% listings and default lstlisting config (remove if unused)
\DeclareNewFloatType{listing}{}
\floatsetup[listing]{style=ruled}

\DeclareCaptionStyle{thesis}{style=base,font={small,sf},labelfont=bf,labelsep=quad}
\captionsetup{style=thesis}
\captionsetup[algorithm]{style=thesis,singlelinecheck=off}
\captionsetup[listing]{style=thesis,singlelinecheck=off}

% Uncomment for table captions on top. This is sometimes recommended by the
% style guide, and even required for some publication types.
%\floatsetup[table]{capposition=top}
%
% (Opinionated rant:) Captions on top are not "compatible" with the general
% guideline that the tables should be formatted to be quickly visually
% comprehensible and *beautiful* in general (like figures), and that the table
% "head" row (with column names) should alone communicate most of the content
% and interpretation of the table. If you just need to show a long boring list
% of numbers (because you have to), either put some effort into showing the
% data in an attractive figure-table, or move the data to an attachment and
% refer to it, so that the boredom does not impact the main text flow.
%
% You can make the top-captions look much less ugly by aligning the widths of
% the caption and the table, with setting `framefit=yes`, as shown below.  This
% additionally requires some extra markup in your {table} environments; see the
% comments in the example table in `ch2.tex` for details.
%\floatsetup[table]{capposition=top,framefit=yes}

\ifEN\floatname{listing}{Listing}
\else\floatname{listing}{Výpis kódu}\fi
\lstset{ % use this to define styling for any other language
  language=C++,
  tabsize=2,
  showstringspaces=false,
  basicstyle=\footnotesize\tt\color{black!75},
  identifierstyle=\bfseries\color{black},
  commentstyle=\color{green!50!black},
  stringstyle=\color{red!50!black},
  keywordstyle=\color{blue!75!black}}

% Czech versions of the used cleveref references (It's not as convenient as in
% English because of declension, cleveref is limited to sg/pl nominative. Use
% plain \ref to dodge that.)
\ifEN\relax\else
\crefname{chapter}{kapitola}{kapitoly}
\Crefname{chapter}{Kapitola}{Kapitoly}
\crefname{section}{sekce}{sekce}
\Crefname{section}{Sekce}{Sekce}
\crefname{subsection}{sekce}{sekce}
\Crefname{subsection}{Sekce}{Sekce}
\crefname{subsubsection}{sekce}{sekce}
\Crefname{subsubsection}{Sekce}{Sekce}
\crefname{figure}{obrázek}{obrázky}
\Crefname{figure}{Obrázek}{Obrázky}
\crefname{table}{tabulka}{tabulky}
\Crefname{table}{Tabulka}{Tabulky}
\crefname{listing}{výpis}{výpisy}
\Crefname{listing}{Výpis}{Výpisy}
\floatname{algorithm}{Algoritmus}
\crefname{algorithm}{algoritmus}{algoritmy}
\Crefname{algorithm}{Algoritmus}{Algoritmy}
\newcommand{\crefpairconjunction}{ a~}
\newcommand{\crefrangeconjunction}{ a~}
\fi


\newenvironment{grammar}[2]
 {\begin{tabular}{@{\qquad}>{$}l<{$}@{\qquad}p{.75\linewidth}@{}}
  \multicolumn{1}{@{}l@{}}{$#1$}&\multicolumn{1}{l@{}}{\hspace{-2em}#2}\\}
 {\end{tabular}}
 % use this file for various custom definitions

\providecommand{\free}[1]{\mtt{free} (#1)}
\providecommand{\dom}[1]{\mtt{Dom} (#1)}
\providecommand{\ran}[1]{\mtt{Ran} (#1)}
\def\cmm{\texttt{C-{}-}}

\newcommand{\calC}{\mathcal{C}}

\newcommand{\li}[1]{\lstinline{#1}}

\newcommand{\mtt}[1]{\mathtt{#1}}
\newcommand{\mbb}[1]{\mathbb{#1}}
\newcommand{\mcal}{\mathcal}

\newcommand{\To}{\Rightarrow}
\newcommand{\Gets}{\Leftarrow}
\newcommand{\GetsTo}{\Leftrightarrow}
\newcommand{\getsTo}{\leftrightarrow}
\newcommand{\pars}[1]{\left(#1\right)}
\newcommand{\pipe}{\ |\ }



\begin{document}

%%% Title page of the thesis and other mandatory pages

%%% Title page of the thesis

\pagestyle{empty}
\hypersetup{pageanchor=false}
\begin{center}

\centerline{\mbox{\includegraphics[width=166mm]{../img/logo-en.pdf}}}

\vspace{-8mm}
\vfill

{\bf\Large MASTER THESIS}

\vfill

{\LARGE\ThesisAuthor}

\vspace{15mm}

{\LARGE\bfseries\ThesisTitle}

\vfill

\Department

\vfill

{
\centerline{\vbox{\halign{\hbox to 0.45\hsize{\hfil #}&\hskip 0.5em\parbox[t]{0.45\hsize}{\raggedright #}\cr
Supervisor of the master thesis:&\Supervisor \cr
\noalign{\vspace{2mm}}
Study programme:&\StudyProgramme \cr
\noalign{\vspace{2mm}}
Study branch:&\StudyBranch \cr
}}}}

\vfill

% Zde doplňte rok
Prague \YearSubmitted

\end{center}

\newpage

%%% Here should be a bound sheet included -- a signed copy of the "master
%%% thesis assignment". This assignment is NOT a part of the electronic
%%% version of the thesis. DO NOT SCAN.

%%% A page with a solemn declaration to the master thesis

\openright
\hypersetup{pageanchor=true}
\pagestyle{plain}
\pagenumbering{roman}
\vglue 0pt plus 1fill

\noindent
I declare that I carried out this master thesis independently, and only with the cited
sources, literature and other professional sources. It has not been used to obtain another
or the same degree.

\medskip\noindent
I understand that my work relates to the rights and obligations under the Act No.~121/2000 Sb.,
the Copyright Act, as amended, in particular the fact that the Charles
University has the right to conclude a license agreement on the use of this
work as a school work pursuant to Section 60 subsection 1 of the Copyright~Act.

\vspace{10mm}

\hbox{\hbox to 0.5\hsize{%
In \hbox to 6em{\dotfill} date \hbox to 6em{\dotfill}
\hss}\hbox to 0.5\hsize{\dotfill\quad}}
\smallskip
\hbox{\hbox to 0.5\hsize{}\hbox to 0.5\hsize{\hfil Author's signature\hfil}}

\vspace{20mm}
\newpage

%%% Dedication

\openright

\noindent
\Dedication

\newpage

%%% Mandatory information page of the thesis

\openright

\vbox to 0.5\vsize{
\setlength\parindent{0mm}
\setlength\parskip{5mm}

Title:
\ThesisTitle

Author:
\ThesisAuthor

\DeptType:
\Department

Supervisor:
\Supervisor, \SupervisorsDepartment

Abstract:
\Abstract

Keywords:
\Keywords

\vss}

\newpage

\openright
\pagestyle{plain}
\pagenumbering{arabic}
\setcounter{page}{1}


\tableofcontents

\chapwithtoc{Introduction}

Systems programming languages address a broad spectrum of challenges that programmers face while developing various kinds of low-level software that typically closely cooperates with hardware, such as operating systems, drivers, real-time mission-critical systems, embedded platforms, and high-performance computation software.
A common systems programming language (such as C) provides a sufficient level of abstraction over the platform to avoid typical errors from repetitive programming in assembly while providing practical guarantees about its code, typically fitting into a highly restricted run-time environment and resource limitations.

Features of such languages have to be carefully selected to reflect these limitations.
The available languages (including C, C++, Rust, Nim, Pascal, and many others \cite{stroustrup1999overview,matsakis2014rust,picheta2017nim,sammet1971brief}) may provide portable structured programming that frees the programmers from the necessity to code assembly directly, a type system that can, e.g., infer types of subexpressions so that they do not need to be declared manually, a static checker (often connected to the type system) that prevents many kinds of programming errors, and facilities for generic programming that promote code reuse.

Here, we loosely follow up on our previous thesis~\cite{klepl2020type} --- we have designed and implemented a variant of a C programming language that supports generic programming driven by a polymorphic type system with typeclass-based overloading.
The system worked, but we have identified several challenges.
Mainly, the type inference for record type members and const-qualifiers would require a type system much stronger than the originally used THIH~\cite{jones1999typing}, and the base language (C) lacked support for more flexible control flow constructions.

The main goal of this thesis is to address these challenges: We extend \cmm, a low-level C-like programming language that was designed as a backend for Haskell compilers, by constructions that allow the programmers to express typeclass-based overloading with multi-parameter typeclasses. That gives the type system sufficient power to infer types of complicated language constructions, including the typed record accesses. To this language, we add a simple typeclass-driven system for automated resource management (similar to resource initialization in C++) and describe and implement an extension of the type system for managing the subtypes (`kinds') specific to \cmm.

As the practical result, create a proof-of-concept compiler of the extended \cmm. The prototype is carried out in Haskell and implements the type inference, monomorphization, several control flow analyses, and a subsequent code generation to LLVM. The thesis additionally provides a comprehensive description of the type inference algorithm based on deferred constraint solving inspired by \citet{vytiniotis2011outsidein}, extended by the subtype constructions required for \cmm. As a partial side result, our subtyping system can be extended to any, even user-specified lattice-based subtype systems, but the subtype inference must be independent of typeclass resolution to avoid complexity issues~\cite{tiuryn1992subtype,frey1997subtype} and we consequently did not pursue a practical implementation.

We show that the compiler prototype can process meaningful programs and output adequate LLVM code and demonstrate the goals of the thesis on small examples. Because the development of the prototype has shown several deficiencies of the compiler and inference algorithm design that were not obvious before later stages of the project, such as the interplay of type inference with monomorphization and subtyping, we also briefly report on the drawbacks introduced by some of the design decisions, and draft several possible future improvements.

\paragraph{Thesis layout.}
The thesis is organized as follows:
\cref{chap1} reviews algorithm W and its suitability for implementing complicated type system features and describes the `French approach' to type inference with deferred constraint solving that allows straightforward inference of multi-parameter typeclasses and existential types.
\cref{chap2} describes \cmm and our extensions to the language for types and resource management.
\cref{chap3} details the design of the compiler and gives a comprehensive description of the type inference algorithm.
\cref{chap4} shows the main results, mainly the functionality of the compiler on selected code examples.
The outcomes and recommendations for future development are briefly summarized in \cref{chap5}, together with several highlights and `lessons learned' from the compiler development process.
We conclude with an overview of thesis goals, achievements, and future work.

\chapter{The French approach to type inference}

\label{chap1}

In this chapter, we discuss possibility of using a modification of the Algorithm W \cite{damas1982principal} for type inference as was the approach of our previous work \cite{klepl2020type}. We then identify the challenges of this approach in the context of type-inferring a C-like language with a Hindley-Milner \cite{damas1982principal} type system extended by multi-parameter typeclasses for overloading. We then introduce an algorithm based on the `French approach', which is simpler to use in the given context and allows for better extensibility.

\section{The Algorithm W and how it connects to deferred solving}

As this thesis builds on the work of the author's bachelor thesis \cite{klepl2020type}, we assume the same syntactic notation and semantics of the untyped lambda calculus, STLC, System F \cite{barendregt1992lambda}, Hindley-Milner type system (HM) \cite{damas1982principal}, and the HM type system with typeclasses described by Typing Haskell In Haskell (here shortened as THIH \cite{jones1999typing}).

We concluded that, if a term is typable in these systems, the appropriate variation of Algorithm W finds it's principal type (with an exception of the System F, type inference of which is undecidable \cite{wells1999typability}.

Algorithm W is an algorithm designed to type-infer the principal types (as described by \citet{damas1982principal}) in a language of lambda terms, outputting the inferred type assignments (a mapping from the input terms to types). This algorithm, in its variation for HM with typeclasses used by THIH, follows the type inference rules listed in \cref{fig:wRules}. From this point forward, if not stated otherwise, by Algorithm W, we understand this variation. Context $\Gamma$ is the set of known type assignments and known predicates (possibly quantified and with further assumptions). Some literature separates this set into multiple ones.

\begin{figure}
    \caption{Type inference rules for the Hindley-Milner type system with typeclasses}
    \label{fig:wRules}
    $$\boxed{\renewcommand{\arraystretch}{2.2}\begin{array}{cc}
        \infer[\text{(Variable)}]{\Gamma \vdash x : \sigma}{x : \sigma \in \Gamma} \\
        \infer[\text{(Application)}]{\Gamma \vdash e e' : {P} \cup {P'}  \Rightarrow \tau}{\Gamma \vdash e : {P} \Rightarrow \tau' \rightarrow \tau & \Gamma \vdash e' : {P'} \Rightarrow \tau'} \\
        \infer[\text{(Abstraction)}]{\Gamma \vdash \lambda x . e : {P} \cup {P'} \Rightarrow \tau' \rightarrow \tau}{\Gamma_x \cup \{x : {P'} \Rightarrow \tau'\} \vdash e : {P} \Rightarrow \tau} \\
        \infer[\text{(Instantiation)}]{\Gamma \vdash e : \sigma'}{\Gamma \vdash e : \sigma & \sigma \geq \sigma'}\\
        \infer[\text{(Generalization)}]{\Gamma \vdash e : \forall \alpha . \sigma}{\Gamma \vdash e : \sigma & \alpha \text{ not free in }\Gamma}\\
        \infer[\text{(Let-polymorphism)}]{\Gamma \vdash \text{let } x = e \text{ in } e' : {P} \cup {P'} \Rightarrow \tau'}{\renewcommand{\arraystretch}{1.1}\begin{array}{c}\Gamma \vdash e : \forall \bar{\alpha} . {P} \Rightarrow \tau \\ \Gamma_x \cup \{x :\forall \bar{\alpha} .  {P} \Rightarrow \tau \} \vdash e' : {P'} \Rightarrow \tau'\end{array}}
    \end{array}}$$
\end{figure}

Throughout this section, we will show that the Algorithm W provides a good basis for typechecking of the type system, but with more extensions, namely multi-parameter typeclasses with functional dependencies \cite{jones2000type} (we abbreviate them MPTCs), it becomes unnecessarily complex and inefficient if we want to use it for type inference as well. We will describe, how the same principles can be used to implement much simpler and more extensible two-phase algorithm based on constraint solving.

Before we begin, let us lay down two terms we will extensively use:

\begin{description}
\item[Type inference] is a process of semantic analysis, which assigns a type (according to the given type system) to each node of the AST, where it makes sense (expressions, procedures, record types, etc.). Each such type can be either monotype or polytype and the object then monomorphic or polymorphic, respectively. Polytypes are characterized, in contrast to monotypes, by being quantified and thus they can represent multiple monotypes.

\item[Monomorphization] statically interprets the generalized callgraph (also including references to type definition) in a type-inferred code and, for every reference to a polymorphic object, it generates a monomorphic copy of this object, which has the polytype appropriately instantiated to the monotype fitting the current context.
\end{description}

\subsection{Limitations of the Algorithm W}

The Algorithm W, which we used in the previous work \cite{klepl2020type}, assumed no recursion in the code and thus any real program with any trivial or nontrivial recursion had to be rewritten into a form where recursive functions were represented by the \lstinline{fix} operator applied to some argument. Type of the \lstinline{fix} operator is $\forall \alpha . \pars {\alpha \to \alpha} \to \alpha$.

The algorithm for preprocessing trivial recursions is derived from fixed point theorem \cite{barendregt1992lambda,damas1982principal}): a function represented as $f \vect{x} = \calC[f, \vect{x}]$, where $C$ is a lambda term such that $\{f\} \cup \vect{x} \subseteq \free \calC$, is simply rewritten into $f = \mathit{fix}\ \lambda F \vect{X} .\calC[f := F, \vect{x} := \vect{X}]$.

For a nontrivial recursion (which we can identify using an algorithm for finding strongly connected components of the call graph). We then use a generalization of fixed-point theorem, so-called double fixed-point theorem \cite{stepanek}, which describes how to construct this rewriting for a component of size two and can be easily generalized for an arbitrary component.

We demonstrated that it is indeed possible to represent the source code in the language of nonrecursive terms, however, doing so is complicated (every language feature has to be mapped to a term) and it does not scale well to possible extensions with more advanced language features (those might introduce further recursive patterns, requiring a redesign of the mapping to terms).

\subsection{Limitations of the syntax-driven aspect of the Algorithm~W}
\label{ex_structs}

In the previous work, for our purposes of extending C to use a HM-style type system with type classes, we determined that the Algorithm W, as described in THIH, is too weak. In the following subsections we explore a possible extension of the Algorithm W and discuss its practicality and how it gives a motivation for the Deferred Solving algorithm introduced in \cref{defer_solve}.

\begin{remark}[Field accessors]
    For the needs of type inference, we, in the context of Algorithm W, represent field accessors as functions from a record structure to the given field (or, equivalently, with a pointer constructor applied to them), each such functions named after the given field name.
\end{remark}

The Algorithm W is able to either type-infer a program with no field overloading (all record field names are unique even between multiple record types) or type-check a program but with limited type inference if we would allow field name overloading.

What does this mean? Let us define two record types:

\begin{lstlisting}
    struct A {int x} a;
    struct B {float x} b;
\end{lstlisting}

Now, if we write \lstinline{a.x} or \lstinline{b.x}, it is obvious that these expressions have types \lstinline{int} and \lstinline{float}, respectively.

If we renamed the field of \lstinline{B} to have different name (for example, \lstinline{y}), we could represent the field accessors as functions $.x : A \to int$ and $.y : B \to float$.

But we do not want to bother the programmer with such strong constraints of requiring the fields to be uniquely named. To allow them have different types even though they share the same name, we have to overload the field accessor functions. Overloading is achieved via introducing a typeclass. We define the field accessors as two instances of a method \li{.x} of ``Has\_x'' typeclass with two type arguments, one for the record type and the other for the field type. The first type argument uniquely represents the instance of the field accessor and the second one allows us to specify the matching field type.

We define the multi-parameter typeclass (MPTC) \li{Has_x}: ($a \to b$ is a functional dependency stating that $a$ uniquely represents $b$)

\begin{center}
    \lstinline/class HasX a b | a -> b {.x : a -> b}/
\end{center}

The two instances are then defined:

\begin{lstlisting}
instance HasX A int {.x : A -> int}
instance HasX B float {.x : B -> float}
\end{lstlisting}

In the HM type system with type classes, this constraint cannot be expressed at all as the class constraints can be defined solely in the terms of taking a single type variable.

In he THIH variation \cite{jones1999typing}, this can be indeed expressed (it is not reflected in the paper, but the supplementary implementation allows arbitrary number of arguments), and the algorithm is strong enough to type-check such construction for concrete types, but it lacks the ability to infer the obvious types of the previously mentioned expressions \lstinline{a.x} and \lstinline{b.x} even though we know the types of both \lstinline{a} and \lstinline{b} and the expressions are defined by them.

\subsection{Tackling the limitations of the Algorithm W}

In \cref{ex_structs}, the inability to type-infer the types of \lstinline{a.x} and \lstinline{b.x} could be possibly solved during the subsequent monomorphization phase, which would know that the type of a record type (\lstinline{a} or \lstinline{b}) determines the types of its fields and would (via lookup in a mapping constructed from their definitions) fill in the appropriate types. In a general case, this would mean that the subsequent monomorphization phase would have to be unnecessarily complex and it would function as a second type inference - consider that solving a type field would allow for solving more constraints and then recursively. So we should address this issue in the type inference itself, in the Algorithm W.

\subsubsection{Using the Algorithm W in the type system with MPTCs}

If we want to simply modify the Algorithm W so it solves principal types even in the presence of MPTCs, we can easily observe that in order to function properly, it has to, on each subprogram (strongly connected component of the program's call-graph; we perform type inference on these subprograms in topological order, beginning with the leaf components), first perform the usual unifications on the syntactic structure of the subprogram, returning a type assignment in the form: $\mtt{name} :: \forall \vect \alpha . P  \To t$, where $P$ is a set of constraints, each having a form $C \vect \tau$, where $C$ is a predicate on the types $\vect \tau$, propagated from the subexpressions of $\mtt {name}$, under some $\vect \alpha = \free t$. And then it has to solve the constraints in $P$ according to some extended constraint solving mechanism for MPTCs, for example with maps, as explored by \cite{jones2000type}. We assume this approach. Here we see the first hints of the modified Algorithm W turning into more complicated version of Deferred Solving we introduce in \cref{defer_solve} as this constraint solving becomes progressively more significant part of the algorithm.

\subsubsection{Intermission: solving class constraints from known proofs}

We show how to solve a class constraint by instantiating a know predicate proof on a small example: given $\Gamma = \{\forall a . \mtt {Ord} [a]\}$, we can prove $\Gamma \vdash \mtt {Ord} [Int]$, using substitution $[a := \mtt {Int}]$. \cite{jones1999typing}

Notice that the substitution $\sigma$ that is used to prove a constraint $C \tau$ from a known predicate proof $C \tau'$ has to satisfy $\tau = \tau' \sigma$, which then also means, it satisfies $\dom \sigma \subseteq \free \tau' \land \ran \sigma \subseteq \free \tau$ (it consumes the free variables of $\tau'$ not present in $\tau$).

Using a more complicated example, we can demonstrate that class constraints often generate additional (more complex) class constraints: To prove $\Gamma \vdash \mtt {Ord} [Int]$ in the context $\Gamma = \{\mtt {Ord}\ \mtt {Int}; \forall a . \mtt {Ord}\ a \To \mtt {Ord} [a]\}$, we first deduce $\Gamma \vdash \mtt {Ord} [Int] \Leftarrow \Gamma \vdash \mtt {Ord}\ Int$ (proven using substitution $[a := \mtt {Int}]$ and the proof $\forall a . \mtt {Ord}\ a \To \mtt {Ord} [a]$) and then $\Gamma \vdash \mtt {Ord}\ Int$ (proven by being an assumption of the context; therefore, using an empty substitution).

Unfortunately, in general, the termination of the process depended on the type system choosing the correct proof --- naively, it could have explored a valid branch that proves $\mtt Ord [a]$ from a proof assuming $\mtt Ord [[a]]$, etc. This could continue indefinitely. This shows that the type system can be pushed to generate infinite chains of proofs for some contexts (if we introduce language extensions such as MPTCs and non-reducing proofs without some appropriate requirements).

\subsection{The challenges in the modified W, a two-phase algorithm}

In the previous section, we showed that each constraint can generate more constraints, and the Algorithm W then solves each of these chains separately. The chains break into trees if one proof has two assumptions, for each of these, then possibly continuing with more assumptions.

The modified Algorithm W with recursion being present in the language alongside functional dependencies we deduced from the record types, is progressively more complicated with each extension, and it should be obvious from the previous examples that the modified Algorithm W is now better described as a two-phase algorithm where the first phase generates some constraints and then the second one solves the constraints by instantiating proofs.

Using the Algorithm W requires some nontrivial preprocessing and mapping the subprogram to the resulting type (as demonstrated by \citet{jones1999typing}), and then, the subsequent constraint solving performs the very inefficient chains of nested proof lookups, as shown in the previous subsection.

We will describe a two-phase algorithm, which does not require the described difficult preprocessing, and introduces ``so-called'' witnesses, which simplify solving the multiple class constraints.

\section{Deferred Solving}
\label{defer_solve}

The idea of using Deferred solving (proposed by \citet{vytiniotis2011outsidein}) is first ``elaborating'' the source program with special variables (elaborations), then running the inference on constraints generated from a semantic analysis of the source program, over the elaboration variables. This allows the constraint solver be independent on the syntactic specifics of the source program as the type semantics of the program are all captured in a simpler \emph{constraint language}. After the inference successfully finishes, we insert the inferred information back into the source program in place of the elaborations.

Elaborations consist of type variables, and then \emph{dictionaries} and \emph{coercions} (defined by \citet{vytiniotis2011outsidein}) that help the subsequent monomorphization phase decide the concrete instances that have to be used at each referring site. In the Deferred Solving algorithm they serve as proofs for each constraint. \emph{Dictionaries} represent type class instances that prove class constraints. And \emph{coercions} prove type identities. We call dictionaries and coercions, collectively, \emph{witnesses}.

Deferred Solving, or Deferred Inference, is very similar to the idea of modified Algorithm W we presented earlier, but it does not require each of the objects appearing in the type-inferred subprogram be directly syntactically mapped to a certain type. It just requires a set of constraints that capture the semantics of the program. These constraints, in the basic version, consist of type equalities between types that have to be equal and class constraints put on certain types that require a proof of the given tuple of types belonging to the given class (these proofs represent instances defined in the type-inferred language).

It also simplifies the notion of type assignments in such a way that all types considered by the algorithm are monotypes. And the type quantification with predicates, which, in the Algorithm W is considered a part of the type assignment, is, in this algorithm, instead, applied to a set of constraints.

We use the following constructs:


\begin{description}
    \item[Type Quantification] A quantification that universally binds a set of type variables.

    We will often omit the word `type', when referring to type quantification.

    \item[Qualification] A construct consisting of a set of predicates (or constraints) being applied to an unquantified <term> of the type language. We then call it a qualified <term>.

    The <term> being qualified is usually a type or a set of constraints.

    \item[Scheme] A <term> scheme is a quantification applied to a qualified term of the type language.

    We will call \emph{type schemes} just schemes.
\end{description}

\subsection{Constraint language}

The language of constraints is an extension of traditional type language with type variables annotated by ``context levels'', giving a language of flat and nested constraints:

\begin{defn}[Language of constraints]
    \label{def:defer_constr}

    Language $F$ of \emph{flat constraints} is defined as follows:

    \begin{center}\begin{grammar}
      \firstcasesubtil{$F$}{d: C \tau^\ast}{Class constraint}
      \otherform{g : \tau_1 \sim \tau_2}{Equality constraint}
    \end{grammar}\end{center}

    The language $W$ of \emph{nested constraints} is defined inductively as follows:

    \begin{center}\begin{grammar}
      \firstcasesubtil{$W$}{F}{Flat constraint}
      \otherform{\forall \alpha^\ast . F^\ast \To W^\ast}{Implication constraint}
    \end{grammar}\end{center}

    \label{constraint_language}
\end{defn}

We call $d$ and $g$ the \emph{witnesses} of the corresponding constraint, because they represent the instances that prove the given constraints.

\begin{remark}
    Unless we allow overlapping proofs \footnote{\emph{overlapping proofs}: proofs such that can instantiate into the same type; equivalently: proofs that have a unification}, the witnesses are isomorphic to the rest of their respective constraint. For example, given two constraints $d_1: C \vect \tau_1$, $d_2: C \vect \tau_1$, then $d_1 = d_2 \GetsTo \tau_1 = \tau_2$.
\end{remark}

Additionally, to the already described equality constraints and class constraints, the constraint language also contains ``implication constraints'', which are constraint schemes. They represent typings, both explicit and implicit, of functions declared by the programmer. Existentials contained in function types (e.g., parameters with run-time type binding, as described in the following subsection) are represented by their own implication constraint.

\begin{defn}[Naming]
    \begin{itemize}
        \item We call the set of constraints, which is then qualified, in the scheme of the implication constraint \emph{nested} in this implication constraint.

        \item We call the set of nested constraints of all the parent implication constraints a \emph{scope} or a \emph{context}.

        \item We use scopes to describe a context level of a constraint. For a constraint, a \emph{context level} is the maximum length of implications leading to it.

        \item \emph{Context level} of the unbound type variables introduced in some scope in \cref{constraint_language} have the context level which is equal to the context level of the scope. \label{constness_level}

        All the appearances of a certain type variable have the same context level.

        \item \emph{Skolem constants} in the given scope are the type variables bound by some parent implication constraint.
    \end{itemize}
\end{defn}

For example, in the following set of constraints
\[ \left\{ a \sim b; \forall c. d \sim e \To \left\{ a \sim d \to e \right\} \right\} \]
the type variables $a, b$ have context level $0$; the type variables $d, e$ have context level $1$; and $c$, being a Skolem constant, has no context level.

\subsection{Existentials}

\label{existentials}

Existentials are types with existentially-quantified type variables, the semantics of such are that there exists some assignment to the existentially-qualified variables that, within the given constraints, describes the value of the existential. Usually, this assignment can be solved only during the dynamic interpretation of the code (i.e., `at runtime').

Semantically, the `universals' (universally quantified types) and existentials differ by that universals have unknown monotypes during definitions and they have to be defined polymorphically, while existentials have unknown types in references and can be defined monomorphically. This leads to an interesting practical consequence: Code specialization required for universal types often leads to code multiplication. On the other hand, use of existentials leads to a necessity of (potentially very unpredictable) use of runtime type and class witnesses, and coercions.

Existentials are typically constrained by some class, which then means that they become witnesses for a method lookup in that class. This can, among other, supply the rest of the code with a type witness that can be derived from the class witness, and does not need to be carried explicitly.

\subsubsection{Existentials as constraints}

Before we start describing the construction of the type system based on constraint solving, we illustrate the aims using a translation of existentials to the language of constraints on an example adapted from \citet{peytonjones2019type}:
The Haskell code in \cref{cex:existentials} shows a definition of an type \lstinline{T} with existentially-typed field \lstinline{getT}, and then a list of variables of that type along with a function that consumes them, accessing the existential.

\begin{codex}
\lstinputlisting[language=Haskell]{examples/existentials.hs}
\caption{Example use of existentials in a Haskell program}
\label{cex:existentials}
\end{codex}

We would like to translate this example into the constraint language. We will decorate the variables in the elaborations with $@$ for clarity, and ignore the elaboration levels for brevity (all global variables are in level 0, and all locally-mentioned variables are in level 1). The functions and their types get elaborated as shown in \cref{cex:exist-elab}, by putting a type variable (``hole'') to each place where the input from the type inference is expected.

\begin{codex}
\begin{lstlisting}[language=Haskell]
MkT :: forall @x . (@y : Show @x) => @x -> T

show :: forall @s . (@r : Show @s) => @s -> String

xs = [ MkT @a (@e : Show @a) 5
     , MkT @b (@f : Show @b) "String"]

f :: @s -> @t
f (MkT @c (@g : Show @c) {getT=(x : @c)} : T) =
    show @d @h x
\end{lstlisting}
\caption{Elaborated program from \cref{cex:existentials}.}
\label{cex:exist-elab}
\end{codex}

The observations from the code are collected into the set of constraints, which is shown in \cref{cex:exist-constr}.

\begin{codex}
\begin{lstlisting}[language=Haskell]
-- from xs:
@a ~ Int,    @e : Show @a,
@b ~ String, @f : Show @a

-- from f:
@s ~ T -- because the input parameter is T
forall @c . (@g : Show @c) => -- w. required by MkT
    @d ~ @c, -- from the type signature of `show'
    @h : Show @d, -- witness required by show
    @t ~ String -- returned by show
\end{lstlisting}
\caption{Constraints generated from \cref{cex:existentials}.}
\label{cex:exist-constr}
\end{codex}

The constraints can be solved by the following assignments $@a := \mtt{Int}, @b := \mtt{String}, @e := \mtt{Show}\ \mtt{Int}, @b : \mtt{Show}\ \mtt{String}, @s := t, @d := @c, @h := @g, @t := \mtt{String}$. The resulting type scheme for \li{f} is: \li{T -> String}.

\FloatBarrier
\subsection{Deferred Inference algorithm}

The Deferred Inference algorithm performs the type inference by rewriting a set of constraints that was generated during the code elaboration into an equivalent, simpler set of constraints in each step, until it either solves all constraints or ends up with a set of constraints that can not be reduced further. \cite{peytonjones2019type}

Our simplified presentation of the algoritm (\Cref{deferred_loop}) is further subdivided to a special case that handles a class constraint (\Cref{case_class}) and an equality constraint (\Cref{case_equality}).

\begin{algorithm}
    \caption{Main loop of Deferred Inference algorithm (simplified, from \citet{peytonjones2019type})}
    \label{deferred_loop}
    \begin{algorithmic}
        \Require $\mtt {Ctx} \equiv \cdots \forall \vect \alpha_2 . A_2 \To \forall \vect \alpha_1 . A_1 \To \mcal W$ is the current context (possibly without the quantifications and corresponding qualifications; or with more)
        \State \Comment $A_1, A_2, \cdots$ are assumptions of the current context $\mtt{Ctx}$
        \State $\mtt{Level} \gets$ the number of quantifications in $\mtt{Ctx}$
        \State $P \gets \emptyset$ is a list of postponed constraints
        \State $\mtt{Continue} \gets \bot$ is a flag signifying whether to continue with postponed constraints
        \While{$\mcal W$ is not empty} \Comment there are unsolved constraints
            \If{$\mcal W = P$} \Comment all constraints are postponed
                \If {$\mtt {Continue} = \top$} \Comment continue with postponed constraints
                    \State $P \gets \emptyset$
                    \State $\mtt {Continue} = \bot$
                \Else
                    \State \Return $\mcal W$ \Comment residual set of constraints
                \EndIf
            \EndIf
            \State $c \gets $ choose from $\mcal W \setminus P$ (according to an arbitrary heuristic)
            \If{$c \equiv d : C \vect t$ is a class constraint}
                \State perform \cref{case_class}
            \ElsIf{$c \equiv g : a \sim b$ is an equality constraint}
                \State perform \cref{case_equality}
            \ElsIf{$c \equiv \forall \vect \beta . P \To \vect W$ is an implication constraint}
                \If{$\vect W$ is empty}
                    \State $\mcal W \gets \mcal W \setminus \{c\}$ \Comment remove the redundant constraint $c$
                \Else
                    \State perform \cref{deferred_loop} with the context $\forall \vect \alpha . A_1 \To \forall \beta . P \To \vect W$, take the resulting $\vect W'$ and update context accordingly
                    \If {a failure $F$ is reported}
                        \State \textbf{propagate failure} $F$
                        \State \Return $\mcal W$
                    \EndIf
                    \State $\mtt{Continue} \gets \mtt{Continue} \lor (\vect W \neq \vect W')$
                \EndIf
            \EndIf
        \EndWhile
    \end{algorithmic}
\end{algorithm}
\begin{algorithm}
    \caption{Solving a class constraint case in Deferred Inference}
    \label{case_class}
    \begin{algorithmic}
        \Require {$c \equiv d : C \vect t$}
        \If{$d$ is a witness of a concrete assumption or proof}
            \State $\mcal W \gets \mcal W \setminus \{c\}$ \Comment remove the redundant constraint $c$
        \ElsIf{$C \vect t$ is proven by some assumption $d' : C \vect t'$}
            \State unify $d$ and $d'$ and globally rewrite all constraints accordingly
            \State $\mcal W \gets \mcal W \setminus \{c\}; \mtt{Continue} \gets \top$ \Comment Solved the constraint $c$
        \ElsIf{$C \vect t$ matches some proof scheme $s \equiv \forall \vect \alpha . A \To d' : C \vect t'$}
            \State instantiate $s$ into $s' \equiv A' \To d' : C \vect t$
            \State $\mcal{W} \gets \mcal{W} \cup A'$ \Comment add $A'$ to the current context
            \State unify $d$ and $d'$ and globally rewrite all constraints accordingly
            \State $\mcal W \gets \mcal W \setminus \{c\}; \mtt{Continue} \gets \top$ \Comment Solved the constraint $c$
        \Else
            \State $P \gets P \cup \{c\}$ \Comment the constraint $c$ is postponed
        \EndIf
    \end{algorithmic}
\end{algorithm}

\begin{algorithm}
    \caption{Solving an equality constraint in Deferred Inference}
    \label{case_equality}
    \begin{algorithmic}
        \Require {$\mtt {Ctx} \equiv \cdots \forall \vect \alpha_2 . A_2 \To \forall \vect \alpha_1 . A_1 \To \mcal W$ is the current context (possibly without the quantifications and corresponding qualifications; or with more)}
        \Require {$c \equiv g : a \sim b$}
        \If{$a$ and $b$ are the same}
        \State $\mcal W \gets \mcal W \setminus \{c\}$ \Comment remove the redundant constraint $c$
        \ElsIf{$\mtt{Level} \geq 1 \land \vect \alpha_1 \cap (\free a \cup \free b) = \emptyset$}
            \State ``promote'' the unification variables of the current context appearing in $c$, globally decreasing their context level, and ``float'' the constraint $c$, moving it into the enclosing (caller's) context; change all contexts (this and the contexts of enclosing loops) accordingly
        \ElsIf{Neither of $a$ and $b$ is a unification variable}
            \If{$a \equiv T \vect s, b \equiv T \vect t, |\vect s| = |\vect t|$} \Comment $a$ and $b$ contain the same type constructor $T$ at the top
                \State $\mcal W \gets \mcal W \setminus \{c\}$ \Comment remove the constraint $c$
                \State $\mcal W \gets \mcal W \cup \{s_n \sim t_n | 1 \leq n \leq |s|\}$ \Comment more equality constraints with corresponding immediate subterms of $a$ and $b$
            \ElsIf{$a \equiv S \vect s, b \equiv T \vect t, S \neq T \lor |\vect s| \neq |\vect t|$} \Comment {$a$ and $b$ contain different type constructors at the top}
                \State \textbf{report unification failure} of $c$
                \State \Return $\mcal W$
            \Else \Comment{Two skolems, type functions, etc.}
            \State $P \gets P \cup \{c\}$ \Comment the constraint $c$ is postponed
            \EndIf
        \Else
            \ {W. L. O. G. Let us assume that $a$ is a unification variable and, if $b$ is one as well, $a$ has an context level greater or equal to $b$'s.}
            \If {$a \in \free b$ }\Comment{Occurs check}
                \State \textbf{report occurs failure} of $c$
                \State \Return $\mcal W$
            \ElsIf {$\mtt{Level} (a) = \mtt {Level}$}
                \State unify $a$ and $b$ and globally rewrite all constraints accordingly
                \State $\mcal W \gets \mcal W \setminus \{c\}; \mtt{Continue} \gets \top$ \Comment Solved the constraint $c$
            \Else
                \State $P \gets P \cup \{c\}$ \Comment the constraint $c$ is postponed
            \EndIf
        \EndIf
    \end{algorithmic}
\end{algorithm}

\citet{vytiniotis2011outsidein} proved that the algorithm is terminating, given some requirements are met in the set of the allowed input constraints. At the same time, the algorithm is confluent --- it produces the same result regardless of reordering of the input.

The precise requirements on the input constraints are rather complicated, we therefore highlight only the ones that are directly relevant to our results:

\begin{enumerate}
    \item All assumptions for witness proofs have to be of the form $d : C a$, where $C a$ is the assumed constraint, $d$ is its witness, and $a$ is a type variable. \label{triv_assump}
    \item All type variables appearing in assumptions of witness proofs have to be free in the proven constraint \label{free_free}
    \item Witness proofs have to be non-overlapping (equivalently: the proven constraints of two distinct proofs cannot be unifiable). \label{noover}
    \item Witness proof schemes have to be of the form $\forall \vect \alpha . P \vect \alpha \To d : C (T \vect \alpha)$, where $P \vect \alpha$ is a set of assumptions, $C (T \vect \alpha)$ is the proven constraint, and $T$ is a type constructor and $d$ is the witness \label{lookup_t}
    \item Ambiguous schemes are not allowed (universally) \label{ambi_rule}
\end{enumerate}

\begin{remark}[Ambiguous assignments]
    We consider an assignment $\mtt{name} :: \forall \vect \alpha . P  \To t$ ambiguous if, after solving all solvable constraints, $\free {P} \setminus \free t \setminus \free \Gamma \neq \emptyset$ (in THIH, there is so-called defaulting described, which can resolve these ambiguities under some strict conditions, but we will not consider them here as they do not apply to general cases).
\end{remark}

The rules \ref{triv_assump}, \ref{free_free} and \ref{lookup_t} together ensure that the steps solving a class constraint are always reducing. The rule \ref{noover} then ensures the class constraints satisfy the confluence property.

The rule \ref{lookup_t} also ensures the proof schemes can be looked up efficiently, as the lookup can be implemented via a lookup table indexed by the $T$ type constructor.

Finally, the rules \ref{triv_assump} and \ref{lookup_t} can also be intuitively generalized for use with multi-parameter type classes.

\subsection{Improvements over the Algorithm W}

\xxx{Tohle potrebuje nejakou uvodni vetu}

\begin{itemize}
    \item  The solving of constraints can be performed in any arbitrary order as long as any possible algorithm extensions preserve its confluence property.

    \item The solving is completely independent on the specifics of the typed language structure, this is achieved via the elaboration it with placeholders for types and witnesses, then expressing the input for inference via those placeholders alone, and then finally filling in the placeholders back to the language.

    \item The notion of witnesses allows us to reduce the amount of explorable proof chains, because the constraints of related or equivalent proofs can be easily unified and removed.
\end{itemize}


\section{Deferred Solving for systems programming}
\label{sys_defer}

In this thesis, we use an algorithm based on the Deferred Inference algorithm because it provides a very easy to extend formalisms which require only small number of modifications to it to achieve the coverage of the desired set of typing capabilities. Namely, MPTCs.

In our type system, we do not support existentials, as we want the maximal transparency of the resulting code for possible use in systems programming. Existentials generate irremovable witnesses as we explained in \cref{defer_solve}, because the actual type of the input existentials is known only in run-time (\citet{grossman2002existential} provides a closer exploration of this topic). This approach guarantees type-safety as it ensures the witnesses are passed properly, but it does so at a cost we would like to not have.

Further, we extended the type system with a highly specialized version of subtyping, which can provide us the necessary means of type-checking and type-inferring whether a variable is semantically a compile-time constant and also what is its `data kind' (we introduce those in \cref{chap2}, which describes the language specifics).

Then, we also introduce functional dependencies. These are implemented via a rewriting scheme that takes precedence over the class constraint resolution. We specify the details in \cref{chap3}.

Not using existentials in our type system allows us to completely omit the context levels from the deferred solving algorithm, which then allows us to solve all equality constraints eagerly, as the solvability of such constraints is limited only by the context levels.

We also solve all constraints of the whole program globally. This turned out to be a working, but very inefficient approach later \xxx{add to conclusion if it is not there}.

\chapter{\cmm{} extension with type inference and automatic resource management}

\label{chap2}

\section{\cmm{} history and significance}

Before we talk about the specifics of \cmm{}, let us explain our choice of this particular language for our demonstration of using type inference and automatic resource management features in a systems-programming context.

Cmm (a variant of \cmm{}) is used as a backend intermediate representation for GHC (see \cite{haskellbackends}), the most popular haskell compiler. \cmm{} is designed on a great experiential basis by Simon Peyton Jones, one of the co-creators of the haskell language and GHC itself.

\begin{remark}
    The \cmm{} language has two major versions, which differ quite significantly.

    It is important to note that, by referring to just \cmm{}, we will always refer to the second version of the language, \cmm{} 2.0 (as specified by \cite{ramsey2005c}).
\end{remark}

Just like Cmm, \cmm{} is designed to serve as a back-end representation for various front-end compilers. It abstracts away various architectural specifics, while it also offers the users to target for some specific architectural features, like specific register sizes, for example, \lstinline{bits32} - a 32 bits long register. It also allows the user to specifically request a tail call instead of a regular function call, which does not constitute stack manipulation.

The register type by itself does not specify whether the register is an unsigned integer, signed integer, floating-point number, pointer, etc. These all are considered as equivalent types in the context of \cmm{}.

One of the main motivations for choosing this language was that it strives to be minimalistic compared to high-level languages while very robust and general, compared to assemblers, offering many features, which can be tied to type inference. It is due to the above reasons, it is usually categorized as pseudo-assembler.

The algorithm our experimental approach is based on, so called ``Deferred Inference'' or the ``French Approach'' (described in the \cref{defer_solve}), was proposed to the general public by the same person, so there were even more arguments for this choice motivated by some overlaps in design philosophy and ideas.

\section{Syntax and properties of \cmm{}}

What characterizes \cmm{} perhaps the best is simply: minimalism. The general syntax for statements and expressions is almost directly a sharp subset of any other C-like language, it supports only bit-vector types with no distinction between integers and float numbers, and the only three control flow statements (not counting call statements) it supports are if statements, switch statements and branches. But, at the same time, it offers a high number of additional compiler hints and constraints the user can specify. All of these characteristics are desirable features of any systems programming language.

\begin{ex}[\cmm{} syntax example]
    \li{sp1} computes the tuple $\pars{\dfrac{n(n+1)}{2}, n!}$ with 32-bit integers
    \begin{lstlisting}
sp1(bits32 n) {
    bits32 s, p;

    if n != 1 {
        s, p = sp1(n - 1);
        return (s + n, p * n);
    }

    return (1, 1);
}
    \end{lstlisting}
\end{ex}

\subsection{\cmm{}: no side-effects in expressions}

What differs the approach of \cmm{} from other C-like languages is that it has call statements instead of call expressions. This is an important distinction, as procedure calls usually perform some hidden side-effects (they have non-pure computation model), hidden in the sense that the side effects would be specified in the language syntax. Not allowing side effects inside expressions is another desirable feature of a systems programming language if we strive for maximum code transparency.

\subsection{Continuations, special labels in \cmm{}}

It is important to note that \cmm{}, while being minimalistic, is not minimalistic in all of its aspects. \cmm{}, in contrast to C and other similar languages, offers so-called ``continuations''. Those are a generalization of label statements that can be called (using a ``cut to'' statement) from one procedure to jump to another. It requires a pointer to the other procedures' activation, and, unlike regular labels, it can be called with values that change the local state of the target procedure: the actual arguments of this ``cut to'' call are bound to local variables inside the activation that is to be executed. Because of this, the fallthroughs to the blocks beginning with these statements are forbidden.

It might seem that these continuations serve as a thread model for the language, but it is not so, they do not implement the whole thread model and the languages does not specify any thread scheduler, but it offers a well-defined hooks for the front-end compilers. Same would apply for any assumptions that this implements some kind of error handling. For more, see \cite{ramsey2005c}

\subsection{Kinds in \cmm{}, an inspiration for a new type system}

As we discussed earlier, \cmm{} does not distinguish between different register types in the type system. The way \cmm{} distinguishes between different register types is by introducing a concept called ``kinds'' (but we will refer to them as ``data kinds'' to distinguish them from kinds we know from the HM type system and similar systems).

Note that these data kinds are usually strictly separated from each other in many languages (by so-called ``aliasing rules'' \xxx{link aliasing rules for c++ (perhaps C)? or for some asssembler if such exist?}).

By specifying a kind, a user can give constraints on which type of physical registers is the value allowed to be stored (or expected to) and then this constraint can be even more tightly constraint by specifying a specific register by its name.

This can be modelled efficiently with a very limited notion of lattice-based subtyping we introduce in the next section, where we consider the empty set of possible registered the bottom and all registers the top(this relates to \cite{tiuryn1992subtype}, which shows \texttt{PSPACE}-completeness for a generalization of this idea).

\section{Extending \cmm{} for type inference}

We extended the language with new syntactic constructs: \li{class}es, \li{instance}s, \li{struct}s, \li{new} specifications, \li{dropped} statements and \li{[ptr]} generic dereference.

And then with \li{auto} anonymous type variables, \li{auto(<name>)} named type variables, \li{ptr(<type>)} pointer types, \li{label} types, and few other types.

And semantically, we extended the language with more intricate type system supporting subtyping, automatic resource management features that define automatic function calls on procedure exits, and class instance (method instance) resolution.

\subsection{Type classes and instances}

\li{class}es and their \li{instance}s are implemented in a similar fashion to the previous work (see \cite{klepl2020type}), and unlike in the previous work, we define functional dependencies.

\li{class} definitions define the user-defined type classes with some superclasses and functional dependencies and then contain list of function declarations defining the methods of the type class. For superclasses and functional dependencies, we borrow Haskell syntax.

Superclasses are requirements that specify that each instance of the class has to correspond to some instances of the superclass; a good example of this would be \li{Eq a} in \li{class Eq a => Ord a} in Haskell, which specifies that all comparable types have to also be equatable.

\begin{ex}[Classes]
    The following code snippet shows a definition of a class \li{C} with an method \li{m}:

    \begin{lstlisting}
        class C a {
            m (auto (a) x) -> auto (a);
        }
    \end{lstlisting}

    The method \li{m} takes an argument of type \li{a} and returns a result of the same type. All instances of this method then have to follow this type scheme.
\end{ex}

\li{instance} definitions then define the instances of their corresponding type classes under some assumptions. For each instance, its assumptions have to satisfy all superclasses of the corresponding class (often, they are the same).

\begin{ex}[Instance]
    The following code snippet then shows a possible instance definition.

    \begin{lstlisting}
        instance C bits32 {
            m (auto x) {
                auto y;
                y = x * 5;

                return (y);
            }
        }
    \end{lstlisting}
\end{ex}

We will often call instances ``proofs'', when discussing them in theory. We require them to be non-overlapping, as, if we allowed them to overlap, the resulting system would be ambiguous (consider, for example, \li{C Int a} and \li{C a Int} in Haskell notation).

\subsection{Structures and member fields}

We added \li{struct}s, a record structure syntax for a feature known from other C-like languages (among others).

Structures use the syntax derived from regular data specialization for maximum consistency. One of the side-effects of this approach is that one field can have various name aliases. This might seem inconsequential, but it can have very interesting consequences in polymorphic programming.

\begin{ex}
    \label{list_ex}
    The following example shows a possible implementation of a common linked list.

    \begin{lstlisting}
        struct list a {
            next: ptr(list auto(a));
            value: data: auto(a);
        }
    \end{lstlisting}

    The last field showcases that the list can be used in both, any function that expects it to carry a value, and in any function that expects it to store some data.
\end{ex}

They allow us to pack more, related, data together (in the \cref{list_ex}, some item and a pointer to the next item) and thus abstract away the specific data layout from the implementations. This then allows us to write functions passing the data or references to it once and then not reimplement them every time the specific data layout changes.

The reason \cmm{} itself does not require structs is that it does not support polymorphism nor it enforces type safety. Instead of passing a pointer to a record and then, in the consumer, using a set of field accessors, it can pass an untyped pointer and then use a set of integer constants representing the offsets of the data.

Since we infer the types and multiple record types can have the fields implemented differently (in the case of our modification, the field might even overlap), the above approach is no longer desirable. It is also very error-prone as the offsets may change with every layout adjustment.

And we also added \li{new} specifiers for \li{struct} objects to implement our automatic resource management goal. For these automatic resource management features see the \cref{RAII}

\subsection{Generic dereference and typed labels}

On the semantic level, we extended the language to have typed labels. All data labels have the type derived from the object stored at the label's location, and the code (branch) labels have a distinct type from all data labels, as they do not point to any such object.

This typing of labels allowed us to make quite subtle, but very useful syntactic extension of allowing the type be omitted in pointer dereference expressions (see \cite{ramsey2005c}), thus introducing a generic dereference.

\begin{ex}[Generic dereference]
    This example showcases that we do not need to specify the type when dereferencing a pointer, the type is automatically inferred.

    \begin{lstlisting}
        stackdata {
            p: bits32;
        }

        [p] = 5; // stores 5 to 'p'
    \end{lstlisting}
\end{ex}

\begin{ex}[Generic dereference with structs]
    This example shows that the generic dereference combines well with structures and their fields.

    \begin{lstlisting}
        struct X {
            x: bits32;
        }

        ...

        stackdata {
            p: X;
        }

        [p->x] = 5; // stores 5 to the 'x' field of 'p'
    \end{lstlisting}
\end{ex}

\subsection{Type system extensions for \cmm{} subtypes}

We extended the language with the HM-like type system with multiple-parameter type classes and functional dependencies. And for this project, we, inspired by the optional data-kind specifications supported by the \cmm{} language, alongside with requirements on various values being computed in compile-time or link-time, then decided to model those in the inference as well. This was very experimental and it provided many insights into intricacies of type systems.

This turned out to be a decision significantly affecting the implementation process as the modelling of these data kinds and constnesses in the context of type inference is regarded as an open question and there were not enough easily accessible materials on primitives that would ensure efficient computability. In this section, we will show the design of such driven by the \cmm{} subtypes as a motivation.

As a result of incrementally building the type system alongside the prototype implementation, the implementation does not currently cover the whole type system space. It required many redesigns and some of its assumptions on the theory had to be either swapped for stronger versions, or completely removed (removed in case of bijections between types and their properties, which we will explore later). The implementation was changed to reflect these refinements in the theory, but it was later concluded, since the theory matured enough, that a complete redesign would be preferable.

\subsubsection{Subtyping system introduction}

We developed a constrained subtyping mechanism extending the HM type system with type classes and type constructors, that type-checks and type-infers what types of registers (integer, float, etc.) a variable can be stored to and when, in terms of run-time, link-time and compile-time, the variable can be assigned a value. This has many implications, amongst which is the possibility of enforcing compile-time evaluation we might know from C++.

This subtyping mechanism models an orthogonal space of ``subtype dimensions'', where each dimension is a lattice with a bottom and we do not consider constraints with sharp inequalities. Sharp inequalities would be very easy way to design an ambiguous system, and also, they do not make any sense when we want to determine just the minimal ``run-timeness'' or the minimal set of registers a value can be received on from an assignment.

We represent the various subtype dimensions of types with so-called ``subtype variables''. We will use these two representations (set of dimensions of a type and a set of variables representing a type's subtyping) synonymously.

\subsubsection{Subtyping system dimensions}

\begin{remark}[Type systems without subtyping]
    We consider the HM type system and all its variations we assumed so far a type system with only one dimension, the typing dimension. It is important to note, that the typing dimension is separate from the type, as having the same typings does not imply the types are necessarily the same.
\end{remark}

\begin{defn}[Subtyping system]
    We define the subtyping system as a minimization problem, where we first minimize the ``typing dimension'' (principal types, see \cite{damas1982principal}) and then the mutually orthogonal ``subtype dimensions''. We introduce these two concepts in definitions \ref{typing_def} and \ref{st_dim}.
\end{defn}

\begin{defn}[Typing dimension]
    \label{typing_def}
    The typing dimension is a lattice with instantiation $\sqsubseteq$ serving as the $\leq_T$ operation (or, equivalently, unification serving as the $\lor$ operation - notice the similarity of concepts: \emph{lowest upper bound} and \emph{most general unification}) and $\forall a . a$ being the bottom element (we assume equivalence under renaming the bound variables, similarly to \cite{barendregt1992lambda}).

    In our proposed type system, we perform the type class resolution using only this dimension. Notice how it is easily specifiable using the language of lattices: $C t \GetsTo \exists C t' \in \Gamma . t \geq_T t'$.
\end{defn}

\begin{remark}
    Note that the unification failure is the top element of the typing lattice. Inferring the unification failure is similar to deriving a contradiction in a logic theory. In both systems, we report such inputs as wrong.
\end{remark}

\xxx{continue from here}

\begin{remark}
    Note that, unlike other, subtype, dimensions (generally; will apply to the ones introduced), the equality in typing dimension is congruent with type applications and type substitutions, and, if we are not considering type functions, all the type operations are bijective.
\end{remark}

\begin{defn}[Subtype dimension]
    \label{st_dim}
    Each subtype dimension $s$ is required to be a lattice with a bottom element, orthogonal to other dimensions, and non-affecting type class instantiation. Again, the top element is allowed to be a failure.

    It is defined by the $\leq_s$ operation (or the $\lor$ operation) and a set of distinct subtyping constants specific to the subtype dimension. This set is nonempty, always containing the bottom $\bot_s$.

    We can then simply extend the definition to apply to user-defined subtype dimensions, but that is not within the scope of this thesis.

    In the type system specific to this thesis, we use two subtype dimensions: \emph{data kinds} and \emph{constnesses}. We will specify them in \xxx{the algorithm}.
\end{defn}


\subsubsection{Subtyping system constraints and variables}

\begin{defn}[Subtype constraints]
    We define the set of constraints as $t \leq_s t'$, where $t$ and $t'$ are two types and $\leq_s$ is ordering of $t$ and $t'$ in the subtype dimension $s$.
\end{defn}

\begin{defn}[Equality of types implies equality of their subtype dimensions]
    If two types $t$ and $t'$ are equal, then $t =_s t$ for every subtype dimension $s$. We specifically do not define this relation as equivalence. This is why we still use the term type to describe the objects and we do not refer just to their typings and subtype dimensions.
\end{defn}

\begin{defn}[(Direct) subtype]
    We call a type $t$ a (direct) subtype of a type $t'$, if $t' \geq_s t$ in each subtype dimension $s$ and also $t =_T t'$ (their typings are the same).
\end{defn}

\begin{defn}[Subtype variable]
    A subtype variable of a type $t$ is a variable $s$, such that $t =_s s$, and, for every other type $t'$, such that $t' \neq t$, and its corresponding subtype variable $s'$, it holds that $t' =_s t \To s = s'$.
\end{defn}

\begin{lemma}
    Lattice-constrained subtyping is efficiently computable at least if every subtype dimension is orthogonal to the mechanism of deciding type class instances.
\end{lemma}

\begin{proof}
    For each topological component of the code (a part of code, that has cyclic references) it is a simple minimization problem on a lattice structure defined by some fixed constants.
\end{proof}

\begin{remark}[Inequality closure convention]
    We say that there is an inequality (or equality) between two (type or subtype) variables if it is specified by a constraint given in the input to the algorithm, it is in a set of base assumptions (relations between subtyping constants, definition of bottom), or comes from transitivity, reflexivity, and (weak) antisymmetry.

    This convention allows us to use more succinct statements when describing the type system.
\end{remark}

\subsection{Properties of the subtype system}
\label{subtyping-idea}

Let us explain the idea of the subtyping system on the following example of inferring the constnesses of a simple procedure.

The procedure's arguments' constnesses are represented by the constness variables $a$ and $b$ and a return value's $c$, and let there be a constant $K$ that represents a concrete constness (perhaps, it can be a link-expr requirement).

If there is a constness variable $x$ coming from some internal variable and then another two constness values $y, w$ and the system contains the following constraints: (we will omit saying more-or-equal in order to make the statements shorter; also, for clarity, it is always better to think in terms ``more constant'' as ``more specific'', rather than, ``more general'')

\begin{itemize}
    \item $x$ is more constant than $y$ (perhaps, because we assign $x$ to $y$. An easy exercise is to show that for these types of subtyping, assignment targets have to be less specific than the source operands)
    \item $y$ is more constant than $c$
    \item $x$ is more constant than $K$
    \item $w$ is more constant than $x$
    \item $a$ and $b$ are more constant than $w$
\end{itemize}

\centerline{\includegraphics[width=\linewidth / 2]{img/out/subtyping-idea.pdf}}

If those are the only subtyping constraints in the given component (we can consider just this subtyping dimension, since they are required to be orthogonal), we can safely unify the constness variables $x$ and $w$ as they are limited by the same fixed values (upper-bounded by $c$ and $K$ and lower-bounded by $a$ and $b$). And we can clearly see that the minimal $x$ will be always equal to $w$ regardless of constnesses of $a$, $b$ and $c$ bound by outside contexts. Less obvious is, whether we can safely unify $x$ and $y$. And since the system is based on minimization, the minimal $y$ will be, again, always equal to $x$ (with an apparent exception of cases where $x$ is not solvable because of $K$; but we are interested in solving the whole program and not single values, so separating them does not make sense - yet still, in practice, this is a good reason for keeping some outside meta-information about the bounds of the original variables for better error messages).

So, to clearly conclude the example. We are interested in an algorithm that unifies variables $x$, $y$ and $w$ and annotates the type of the function with constraints $c \geq a, c \geq b, K \geq a, K \geq b$ for the ``outside subtype information'' and $c \geq x, K \geq x, x \geq a, x \geq b$ for the ``inside subtype information'' (note that the other two variables are not mentioned, as they are unified with $x$).

We then will go even further and state that $x = \sup \{a, b\}$ (and that $x \leq \inf \{c, K\}$), in other words, $x$ is defined by its fixed lower-bounds. Notice the equality in the former constraint. We then can completely remove the latter constraint $x \leq \inf \{c, K\}$ as it is redundant (can be shown via transitivity from the outside subtype information).  We then will use a similar notation for more understandable outside subtype information for each variable. For $c$, that would be $c \geq \sup \{a, b\}$. Fixing the subtype variables to suprema of their fixed lower-bounds has an effect of not having any free variables in the constraints for the procedure.

This example demonstrates that the minimal subtypes within a space of fixed values can be transformed into an equality theory.

\begin{remark}[Sharp inequalities]
    If the theory contained constraints with sharp inequalities, we could not separate the outside subtype information from the inside subtype information, nor have such liberty in unifying the variables.
\end{remark}

\subsection{Inference of subtypes}

We are using the term ``fixed values'' in the \cref{subtyping-idea}, but we have not yet introduced its meaning. They are some constants that are not to be optimized by the subtyping inference mechanism. In the most simple version of this system, if we do not consider type classes (only the HM type system with type constructors), the fixed values consist of subtyping dimension constants (in the \cref{subtyping-idea}, that would be $K$); and then the subtype dimensions of skolem constants (again, in the \cref{subtyping-idea}, represented by the symbols $a$, $b$, and $c$).

\begin{defn}[fixed value in the subtyping system without type classes]
    Each subtype variable of a skolem constant or subtyping dimension constant is a fixed value.
\end{defn}

\begin{defn}[(Closest) fixed lower bounds]
    We will call ``(closest) fixed lower bounds'' of a subtype variable $v$ in a subtype dimension $s$ such fixed values $c$ in the dimension $s$, that $c \leq_s v$ and there is no other $c'$, such that $c \leq_s c' \leq_s v$ would be in the closer of the known inequalities.
\end{defn}

\begin{remark}[Monotypes in the subtyping system]
    If the type is a monotype (and thus it is a valid type for a program), it has no skolem constants and that means we can safely define all the free subtype variables as suprema of their fixed lower bounds.
\end{remark}

\begin{lemma}[Each subtype variable can be set to a suprema of its lower bounds, where each such bound is the first fixed value on each path from this variable towards the bottom; this will give us the minimimal solution to the constraints]
    \label{suprema_subtyping}
\end{lemma}

\begin{proof}
The validity of this simplification can be shown in toplogical order of the variables, where the topology is defined by their inequalities. If they are in a strongly connected component, they have to be equal, so without a loss of generality, we can assume, that there is a topological ordering of the variables, starting with the lowest.

If a set $X$ denotes the set of all the lower bounds of a variable $x$ and it is the variable with the lowest topological order. That means that if a variable $y$ constrained by $x$, then the constraint is $x \leq y$. If $x^\star$ is equal to the supremum of its lower bounds, then for every other $x'$ such that $x' \geq x^\star$, if $y \geq x'$, then also $y \geq x^\star$. This means that $x^\star$ is at least as good as $x'$ for all the constraints regarding $x$, while it also minimizes $x$ itself.
\end{proof}

\begin{defn}[Suprema subtype constructor]
    The suprema subtype constructor $\sup_s$ is defined as a function from a set of values of a subtype dimension $s$ to a value from this dimension. It is defined by the suprema operation on the lattice defining the subtyping dimension.

    Therefore it holds the following identity: $\sup \{v\} = v$. Note that because of this, it is not equivalent to a type constructor. But, just like a type constructor, it defines a value, albeit nonuniquely.
\end{defn}

\begin{defn}[Outside subtype information (of a function)]
    The outside subtype information in a subtyping dimension $s$ consists of subtype constraints between all fixed values in the subtyping dimension $s$. And the bindings to their respective types. All subtype variables appearing in a strongly connected component are unified.

    It can be reduced under transitivity and reflexivity for acquiring a canonical form.
\end{defn}

\begin{remark}[Outside subtype information]
    The bindings of types to their corresponding subtype variables are $1 \to \star$.
\end{remark}

\begin{defn}[Inside subtype information]
    The inside subtype information consists of a set of statements in the form $x = \sup X$ and the bindings to their respective types. It is not a part of the type signature of the function.

    Validity of this is described by the \cref{suprema_subtyping}.
\end{defn}

\begin{remark}[Typechecker for user-code that contains blackbox function does not need the inside subtype information]
    Another direct consequence of the \cref{suprema_subtyping}.
\end{remark}

If we extend this mechanism to the type system with typeclasses, we have to regard the subtype dimensions of variables constrained in a type class constraints.

\begin{defn}[fixed value in subtyping with type classes]
    Each subtype variable of a skolem constant, subtyping dimension constant, or a subtype variable of a type variable appearing in a type class constraint is a fixed value.
\end{defn}


The algorithm we based our type inference on uses context levels that measure how free the variables are in the given context, see the \cref{constness_level}. This applies to their subtype dimensions as well

\begin{defn}[fixed value in subtyping with type classes and context levels]
    Each subtype variable of a skolem constant, subtyping dimension constant, a subtype variable of a type variable appearing in a type class constraint, or a subtype variable with a lower context level than that of the solver (or equivalently, scope) is a fixed value.
\end{defn}

We did not strive for representing existentials, as there has to be runtime type information included with them, so we will not consider such system throughout this thesis. Yet, we specify it for any such extensions.


\subsection{Automated resource management extensions}
\label{RAII}

\xxx{\url{https://www.appinf.com/download/C++BetterDeviceSoftware.pdf}}
\xxx{\url{https://www.stroustrup.com/bs_faq2.html\#finally}}


\xxx{check this}

The RAII (Resource acquisition is initialization) is a mechanism and a design pattern, which can be used for automatic execution of a code (commonly referred to as ``resource release''), which is meant to be run ``in every case''. It is usually used with locks (so-called ``mutexes'') and various ``smart'' pointers and containers.

In the language the term is coined by, is characterized by tying this action (resource) to a certain object's lifetime and the action then, through the reference to the object, usually frees a lock, or deallocates a memory, decrements a counter, etc.

It is directly tied to the ``try - finaly'' \xxx{ref stroustrup} construct, but it requires less code with less nesting in the scope it is used in, and also, it allows for more flexibility. Which we will point out in the conclusion \xxx{ref}.

The statement that this construct requires less code is supported by the claim that there are usually far more resources than the kinds of these resources.

We extend the type system by introducing the \li{new} specifier, which signifies, that an stack object is to be tied to some predefined automatic resource management action. This action is defined by the \li{drop} function, which takes a pointer to the resource-representing object. This function is called with every resource-representing object on every exit point of the procedure.

We also introduce the \li{dropped} statement, which specifies, that the resource is to be considered ``dropped`` on all exit points following this statement.

\begin{ex}[Trivail use of automatic resource management]
    This example shows a common use of a resource-representing object, \li{X} can be, for example, some mutex
    \begin{lstlisting}
stackdata {
    new X;
}

...

// implicitly calls drop(ptr(X) _); on the anonymous object
return();
    \end{lstlisting}
\end{ex}

\begin{remark}[Explicit initialization]
    Note, that the resource initialization is to be called explicitly.

    The reason for this design choice is that the resource release much more often does not require any extra arguments, which would specify, how it should be performed.

    And, at the same time, the resource initialization very often requires a different behavior on each occurrence. Consider, for example opening and closing a file.
\end{remark}

\begin{ex}
    This example then shows more complex use of automatic resource management, where we specify a part of the current scope, where the resource release is to be done explicitly or not at all.

    \begin{lstlisting}
stackdata {
    x: new X;
    // notice that, because of the dropped statement,
    // we named the resource-representing object.
}

if a == 5 {
    dropped x;

    // here, the user can explicitly call drop(x)
    // if they want

    // does not call drop(x);
    return ();
} else {
    // implicitly calls drop(x);
    return();
}
    \end{lstlisting}
\end{ex}

\chapter{Design}

\todo{PROC BC JO ALE PROC ZLEPSIT}

\todo{VRSTVY}

\section{Flattening and blockifying}

\section{Flow analysis and how it connects to RAII}

\section{Type System}

\section{Inference}

\subsection{Kind inference}

\todo{add text from kinds.tex}

\subsection{Type inference}

\begin{defn}[Type variables]
    $\mathbb{V}$ is the space of type variables given as follows, each (sub)type considered by the inference algorithm will be assigned a type variable representing it:
    \begin{grammar}{\mathbb{V} \Rightarrow v}{}
        \mid \mathbb{V} '
    \end{grammar}

    The variables follow the obvious lexical ordering (based on their length).
\end{defn}

\begin{defn}[Typings]
    $\mathbb{T}$ is the space of typings given as follows, they specify that a type follows the same typing as a different type:
    \begin{table}[H]
        \begin{grammar}{\mathbb{T} \Rightarrow }{Typing}
            \mathbb{V} & Variable \\
            \left[\mathbb{T}\right] \to \mathbb{T} & Function \\
            \mid \left[\mathbb{T}\right] & Tuple \\
            \mid \mathbb{T}\ \mathbb{T} & Application \\
            \cdots
        \end{grammar}
    \end{table}
\end{defn}

\begin{defn}[Types]
    Types use the same language as typings and thus the space $\mathbb{T}$ is considered the space of types as well as the space of typings.
\end{defn}


\begin{defn}[Primitive patterns]
    $\mathbb{P}$ is the space of primitive patterns given as follows, they represent the topmost level of a type's structure:
    \begin{table}[H]
        \begin{grammar}{\mathbb{P} \Rightarrow }{Primitive pattern}
            \left[\mathbb{V}\right] \to \mathbb{V} & Function \\
            \mid \left[\mathbb{V}\right] & Tuple \\
            \mid \mathbb{V}\ \mathbb{V} & Application \\
            \cdots
        \end{grammar}
    \end{table}
\end{defn}

\begin{lemma}[Primitive patterns]
    \label{typesObs}
    Every type $\tau \in \mathbb{T}$ can be equally represented by a type variable and a set of primitive patterns $P \in \mathbb{P}$. This also provides us the way to unambiguously identify the type $\tau'$ of a subexpression $x'$ ($x' : \tau')$ of the expression $x$ ($x : \tau$).
\end{lemma}

\begin{defn}[Constnesses]
    $\mathbb{C}$ is the space of constnesses given as follows:
    \begin{table}[H]
        \begin{grammar}{\mathbb{C} \Rightarrow }{Constness}
            \mathtt{Regular} & \textit{for regular runtime variables} \\
            \mid \mathtt{Linkexpr} & \textit{for variables that are to be known during linkage} \\
            \mid \mathtt{Constexpr} & \textit{for variables that are to be known during compilation}
        \end{grammar}
    \end{table}

    Constness of a type specifies the time when an object assigned the type is to be computed.

    We define ordering on constnesses: $\mathtt{Constexpr} < \mathtt{Linkexpr} < \mathtt{Regular}$.
\end{defn}

\begin{defn}[Data kinds]
    $\mathbb{K}$ is the space of data kinds, it is a subset of the power set of a set of registers.

    The kinds are ordered by the relation of being a subset. We put $\mathtt{GenericKind}$ as the superset of all possible kinds and $\mathtt{EmptyKind}$ as the subset of all possible kinds.
\end{defn}

\begin{defn}[Type properties]
    $\mathbb{D}$ is the space of type properties given as follows:
    \begin{algorithmic}
        \State{$\mathbb{D} \Rightarrow (typing: \mathbb{T}, constness: \mathbb{V}, kind: \mathbb{V})$}
    \end{algorithmic}

    This states that a type has a certain typing, constness of a certain type and a kind of a certain type.

    The functions that project objects from $\mathcal{D}$ to the respective properties are called $t$, $c$, $k$.
\end{defn}



\begin{defn}[Facts]
    $\mathbb{F}$ is the space of flat facts given as follows:
    \begin{table}[H]
        \begin{grammar}{\mathbb{F} \Rightarrow }{Flat facts}
            \mathbb{T} \sim \mathbb{T} & Unification: \textit{specifies that the types have to be identical} \\
            \mathbb{T} \sim_T \mathbb{T} & Typing unification: \textit{specifies that the types have to follow the same typing} \\
            \mathbb{T} \leq_K \mathbb{T} & SubKind: \textit{specifies that the right operand has more general data kind than the first operand} \\
            \mathbb{T} \leq_C \mathbb{T} & SubConst: \textit{specifies that the right operand has more general constness than the first operand} \\
            \mathbb{K} \leq_K \mathbb{T} \leq_K \mathbb{K} & Kind bounds: \textit{specifies that the type has a data kind from the given range} \\
            \mathbb{C} \leq_C \mathbb{T} \leq_C \mathbb{C} & Const bounds: \textit{specifies that the type has a constness from the given range} \\
            C\ \mathbb{T} & Class constraint: \textit{specifies that the given constraint type (applied to a certain number of types - unless nullary) is valid} \\
            \mathbb{T} \leq_M \mathbb{T} & Instantiation: \textit{specifies that the left operand is monotype instance of the polytype right operand; note: this is satisfiable only if the right operand is indeed a polytype} \\
        \end{grammar}
    \end{table}

    Sometimes we will use derived facts like $\mathbb{T} \leq_K \mathbb{K}$, $\mathbb{T} =_K \mathbb{T}$, etc. These are either special cases of the aforementioned facts, or compositions of multiple facts. The most complicated derived fact we will use is the ``SubType'' fact, $\mathbb{T} \leq \mathbb{T}$, which is a combination of $\mathbb{T} \sim_T \mathbb{T}$, $\mathbb{T} \leq_K \mathbb{T}$ and $\mathbb{T} \leq_C \mathbb{T}$.

    It should be noted that, within the defined language, we cannot state the facts $\mathbb{T} \not\sim_T \mathbb{T}$, $\mathbb{T} \not\leq_K \mathbb{T}$, $\mathbb{T} \not\leq_C \mathbb{T}$ and $\mathbb{T} \not\leq_M \mathbb{T}$ and thus these (and similar ones) are not derived facts.

    $\mathbb{W}$ is the space of facts given as follows:

    \begin{table}[H]
        \begin{grammar}{\mathbb{W} \Rightarrow }{Facts}
            \mathbb{F} & Fact \\
            \forall \alpha_1, \dots \alpha_n . [\mathbb{F}] \Rightarrow [\mathbb{W}] & Nested facts: \textit{specifies that satisfying the given facts requires the nested facts to be satisfied as well}
        \end{grammar}
    \end{table}
\end{defn}

\begin{defn}[Bounds]
    We call the rightmost operand of a bounds fact the upper-bound and the leftmost operand the lower-bound.
\end{defn}

\begin{defn}[Trivial bounds]
    If a fact $F$ states $k_1 \leq_K \tau \leq_K k_2$ and $k_2 \leq k_1$, or $c_1 \leq_C \tau \leq_C c_2$ and $c_2 \leq c_1$, we call such a bounds fact trivial.
\end{defn}

\begin{lemma}[Trivial bounds]
    \label{trivBoundsObs}
    \begin{enumerate}
        \item If the bounds of a trivial bounds fact are equal, the type constesses of the types constrained by such a fact have to be the same and we can unify them.
        \item If, for a bounds fact, the upperbound is lower than the lowerbound, the fact cannot be satisfied.
    \end{enumerate}
\end{lemma}

\begin{defn}[Algorithm state]
    The algorithm state consists of the following finite variables
    % TODO: explain each
    \begin{itemize}
        \item Facts: $\mathcal{F} :: [\mathbb{W}]$
        \item Active type variables: $\mathcal{V} :: \subset \mathbb{V}$
        \item Forgotten type variables: $\mathcal{G} :: \subset \mathbb{V}$
        \item primitive patterns: $\mathcal{P} :: \subset \mathbb{P}$
        \item Type properties: $\mathcal{D} :: \subset \mathbb{D}$
        \item SubKinds: $\mathcal{K} :: \text{Directional graph}\ \mathbb{V}$
        \item KindBounds: $b_K :: \mathbb{V} \to \text{interval} \left[\mathbb{K}, \mathbb{K}\right]$
        \item SubConsts: $\mathcal{C} :: \text{Directional graph}\ \mathbb{V}$
        \item ConstBounds: $b_C :: \mathbb{V} \to \text{interval} \left[\mathbb{C}, \mathbb{C}\right]$
        \item Type explanation: $p :: \mathcal{V}_p \to \mathcal{P}; \mathcal{V}_p \subseteq \mathcal{V}$
        \item Type definition: $d :: \mathcal{V} \to \mathcal{D}$
        \item Result function: $u :: \mathcal{G} \to \mathcal{V}$
    \end{itemize}

    By $<_{\mathcal{K}}$ and $<_{\mathcal{C}}$,
    we understand that there exist a directional path between the given distinct operands in the respective graphs.
    We naturally extend this definitions to the other comparison operators
    (for example: $x =_{\mathcal{K}} y$ if there are paths back and forth between $x$ and $y$ in the graph $\mathcal{K}$ or $x = y$). % TODO: not that they can be equal

\end{defn}

\begin{defn}[Algorithm intrastate invariants]
    The algorithm is designed to have the following intrastate invariants:

    % TODO: invariant about u

    \begin{enumerate}
        \item Active and forgotten variables: $Var(\mathcal{F}) \cup Var(\mathcal{P}) \cup Var(\mathcal{D}) \cup Var(\mathcal{K}) \cup Var(\mathcal{C}) = Var(\mathcal{V})$ and $Var(\mathcal{V}) \cap Var(\mathcal{G}) = \emptyset$. \label{invVar}

        \item $\mathcal{V}$ and $\mathcal{D}$ are bijected by the type definition function $d$ with its inverse $d^{-1}$. The subset $\mathcal{V}_p \subseteq \mathcal{V}$ and the set $\mathcal{P}$ are bijected by the type explanation function $p$ and its inverse $p^{-1}$. Note that there generally are many types with no known explanation. \label{invPD}

        We expand the function applicability of $b_K$ and $b_C$ to type definitions is such a way that for an arbitrary type variable $v \in \mathcal{V}$, we have $b_K (d (v)) = b_K (c (d (v)))$, and similarly for $b_C$. In other words: applying $b_K$ to a type definition is equivalent to applying it to the type's kind (and similarly for $b_C$).

        We expand the applicability of $p$ and $d$ in such a way that for an arbitrary type variable $v \in \mathcal{V}_p$ it holds that $d(v) = d(p(v))$ and $p(v) = p(d(v))$.

        And finally, we expand the applicability of $t, k, c$ in such a way that for an arbitrary $v \in \mathcal{V}$ it holds that $t (v) = t (d (v))$, $k (v) = k (d (v))$, and $c (v) = c (d (v))$.

        \item 

        \item \label{invG} If we have two type variables $v, v' \in \mathcal{V}$ representing two types, then:
            \begin{enumerate}
                \item If $k (v) =_{\mathcal{K}} k (v')$, then: $k (v) = k(v')$. In other words, graph $\mathcal{K}$ has no strongly connected components (this will be true for $\mathcal{C}$ as well).
                \item If $c (v) =_{\mathcal{C}} c (v')$, then: $c (v) = c(v')$.
                \item If $b_K (k (v)) = b_K (k (v')) = [k_1, k_1]$ for some kind $k_1 \in \mathbb{K}$, then $k (v) = k (v')$.
                \item If $b_K (k (v)) = b_K (k (v')) = [k_1, k_2]$ for some kinds $k_1, k_2 \in \mathbb{K}$, $k_2 < k_1$, then $k (v) = k (v')$ and it is an invalid kind.
                \item If $b_C (c (v)) = b_C (c (v')) = [c_1, c_1]$ for some constness $c_1 \in \mathbb{C}$, then $c (v) = c (v')$.
                \item If $b_C (c (v)) = b_C (c (v')) = [c_1, c_2]$ for some constnesses $c_1, c_2 \in \mathbb{K}$, $c_2 < c_1$, then $c (v) = c (v')$ and it is an invalid constness.
                \item If $k (v) <_{\mathcal{K}} k (v')$, then for bounds $k_1, k_2, k_3, k_4$ such that $b_K (k (v)) = [k_1, k_3]$, and $b_K (k (v')) = [k_2, k_4]$ it holds that $k_1 \leq k_2$ and $k_3 \leq k_4$. Note that for new (and any non-stored) variables we will assume the least restrictive bounds, so adhering to this invariant means we update the bounds monotonically.
                \item If $c (v) <_{\mathcal{C}} c (v')$, then for bounds $c_1, c_2, c_3, c_4$ such that $b_C (c (v)) = [c_1, c_3]$, and $b_C (c (v')) = [c_2, c_4]$ it holds that $c_1 \leq c_2$ and $c_3 \leq c_4$.
            \end{enumerate}

        \item If a type represented by a type variable $v$ is explained ($v \in \mathcal{V}_p$), then $t (d (v)) = p(v) \left[ \tau := t (d (\tau)) | \tau \in \mathrm{free} (p (d))\right]$. In other words, the typing of the type represented by $v$ follows the typing generated from the typings of the types via which it is explained. \label{invT}
    \end{enumerate}
\end{defn}

\begin{defn}[Algorithm interstate invariants]
    The algorithm is designed to have the following interstate invariants: (we use the indices to distinguish a successor state variable from its predecessor, we will always assume $n < m \in \mathbb{N}$, we then use these indices in all derived functions too).

    \begin{enumerate}
        \item Forgotten variables stay forgotten: $v \in \mathcal{G}_n \Rightarrow v \in \mathcal{G}_m$.
        \item Forgotten variables used to be active: $v \in \mathcal{G}_m \Rightarrow \exists n . v \in \mathcal{V}_n$.
        \item All free variables from facts are considered by the algorithm: $v \in \mathrm{free} (\mathcal{F}_0) \Rightarrow v \in \mathcal{V}_0$.
        \item Assumed subtype constraints are monotonically getting restricted. For kinds: $v \in {\mathcal{V}}_n \Rightarrow {b_K}_m (k_m(u_m(v))) \subseteq {b_K}_n (k_n(v))$ and for constnesses: $v \in {\mathcal{V}}_n \Rightarrow {b_C}_m (c_m(u_m(v))) \subseteq {b_C}_n (c_n(v))$.
        \item Typings do not get forgotten: $v \in {\mathcal{V}}_n \Rightarrow \exists s . s(t_n(v)) = t_m(u_m (v))$.
        \item Explanations do not get forgotten: $v \in {\mathcal{V}_p}_n \Rightarrow \exists s . s(p_n(v)) = p_m(u_m (v))$.
        \item Results are results of results: $u_m = u_m (u_n)$.
        % TODO: add more constraints
    \end{enumerate}
\end{defn}

\begin{defn}[Unifications]
    In the scope of this algorithm we distinguish 4 types of unification. The algorithms for each will be the same, but their meaning is different.

    \begin{enumerate}
        \item Type unification: we unify two type variables that represent types, we then apply the resulting unification to all type variables considered by the algorithm. \label{tUni}
        \item Typing unification: we unify two typings of certain types, we then apply the resulting unification to all other typings considered by the algorithm.
        \item Kind unification: we unify two type variables that represent kinds of certain types, we then apply the resulting unification to all other variables representing kinds considered by the algorithm. \label{kUni}
        \item Constness unification: we unify two type variables that represent constnesses of certain types, we then apply the resulting unification to all other variables representing constnesses considered by the algorithm. \label{cUni} 
    \end{enumerate}
\end{defn}

\begin{defn}[Fresh variable]
    \label{freshVar}
    A fresh variable is chosen according to: \linebreak $\inf \left(\mathbb{V} \setminus \mathcal{V} \setminus \mathcal{G}\right)$.
\end{defn}

\begin{defn}[Algorithm middlesteps]
    \label{middlesteps}
    After each step of the algorithm that transforms one of the state variables (mainly $\mathcal{G}$,$\mathcal{P}$ or $\mathcal{D}$) into a new state, we perform the minimal necessary unifications to ``repair'' the bijections $d$ and $p$ (intrastate invariant \ref{invPD}, all unifications), to collapse any strongly connected components of the graphs $\mathcal{K}$ and  $\mathcal{C}$ and to remove any duplicate images of trivial bounds from $b_K$ and $b_C$ (intrastate invariant \ref{invG}, unifications \ref{kUni} and \ref{cUni}). This will be often performed recursively. We then apply the resulting combined unification (limited to just unifications \ref{tUni}) to the function $u$ so it satisfies the intrastate invariant \ref{invVar}. % TODO: fix wording
\end{defn}

\begin{lemma}[Algorithm middlesteps]
    The algorithm may (correctly) fail during the middlesteps. This can be shown by the observation \ref{trivBoundsObs}.
\end{lemma}

\begin{defn}[Introducing new (fresh) variable]
    \label{introVar}
    We introduce a new (fresh) variable by extending the sets $\mathcal{V}$ and $\mathcal{D}$
    and the type definition function $d$ by the argument $v$ and its image $i = (typing: v, constness: v, kind: v)$, formally: $\mathcal{V}' = \mathcal{V} \cup \{v\}$, $\mathcal{D}' = \mathcal{D} \cup \{i\}$, $d' = d [v := i]$. Where $v$ is the new (fresh) variable.
\end{defn}
\begin{defn}[Algorithm presteps]
    \label{presteps}
    All the (flat) facts generally use complicated types, before each step of the algorithm we replace the types in the observed fact with type variables that represent the types.

    The transformation is performed according to the observation \ref{typesObs}, changing the state variables $\mathcal{P}$, $\mathcal{D}$, $\mathcal{V}$, $p$ and $d$ accordingly and we set the typing of the type to adhere to the intrastate invariant \ref{invT}. For the yet unnamed types we perform the steps explained in definitions \ref{freshVar} and \ref{introVar}.
\end{defn}

\begin{defn}[Algorithm initialization]
    The initial state is: $\mathcal{V} = \mathcal{G} = \mathcal{P} = \mathcal{D} = \mathcal{K} = \mathcal{C} = \emptyset$, $p = d = u = b_K = b_C = \emptyset \to \emptyset$ and $\mathcal{F}$ is the list of facts to be satisfied.

    Before the algorithm starts, we introduce the variables encountered in $\mathcal{F}$ according to the definition \ref{introVar}.
\end{defn}

\begin{defn}[Algorithm steps]
    After the prestep (definition \ref{presteps}), we retrieve the first fact $f$ of $\mathcal{F}$ and then:

    \begin{itemize}
        \item If $f \equiv t_1 \sim t_2$, we unify $t_1$ and $t_2$
        \item If $f \equiv t_1 \sim_T t_2$, we unify the typings $t (d (t_1))$ and $t (d (t_2))$, resulting in the unification $m$ which we then apply to all typings in $\mathcal{D}$.
        \item If $f \equiv t_1 \leq_K t_2$, we add an edge $(t_1, t_2)$ to $\mathcal{K}$.
        \item If $f \equiv t_1 \leq_C t_2$, we add an edge $(t_1, t_2)$ to $\mathcal{C}$.
    \end{itemize}

    Then we perform the middlestep (definition \ref{middlesteps}) and continue to the next step.
\end{defn}

\section{Monomorphization}

\section{Translation}

\chapter{Results and discussion}
\label{chap4}

In this chapter, we show the possibilities of use of the extended \cmm.

\section{Prototype compiler of typed \cmm}
\label{sec:protoComp}

\subsection{Type-inferring local declarations}

First, we demonstrate simple properties of the type system. As expected, the compiler is able to infer types of local declarations. On a simple example of \li{gcd} procedure (\cref{lst:localDecl}) that contains variables declared with ``auto'' type, it is able to deduce the correct type for a local variable and the return value.

\begin{listing}
    \caption{Example of type-inferring local declarations}
    \label{lst:localDecl}
    \begin{center}
    %\begin{minipage}[t]{0.5\linewidth}
    \begin{minipage}{0.5\linewidth}
    \begin{lstlisting}
// input code
foreign "C" gcd(bits32 u,
                bits32 w)
                -> auto {
  auto r;
  while:
  if w != 0 {
    r = w;
    w = u % w;
    u = r;
    goto while;
  }

  return (u);
}
    \end{lstlisting}
    \end{minipage}%
    \begin{minipage}{0.5\linewidth}
    \begin{lstlisting}
// type-inferred code
foreign "C" gcd(bits32 u,
                bits32 w)
                -> bits32 {
  bits32 r;
  while:
  if w != 0 {
    r = w;
    w = u % w;
    u = r;
    goto while;
  }

  return (u);
}
    \end{lstlisting}
    \end{minipage}
    \end{center}
\end{listing}

\subsection{Reduction of code using polymorphism}
\label{sec:reduction}

Type inference allows us to write polymorphic procedures, which are then correctly interpreted by the compiler and rewritten into an efficient code with correctly filled types for each each use in different context. The user does not have to define a procedure for each use separately. This is demonstrated by \cref{lst:reduction}. The code is ``monomorphized'' into code containing monomorphic definitions replacing the user-provided polymorphic ones with correctly inferred types.

Notably, compared to some other solutions (for example: Java, Python), which provide various ways to reduce code by polymorphism, the resulting code contains no runtime checks of types and the code is as efficient as if it was written by hand for every specific type.

This constitutes ``mangling'' the names of monomorphic copies of polymorphic procedures depending on their types. This is performed so each generated definition has a unique name and they can be correctly referenced from their corresponding call sites.

\begin{listing}
    \caption{Example of reduction of code (some names truncated for brevity)}
    \label{lst:reduction}
    \begin{center}
    %\begin{minipage}[t]{0.5\linewidth}
    \begin{minipage}{0.5\linewidth}
    \begin{lstlisting}
// Shorter input code
// User-provided
// polymorphic definition
add(auto x, auto y) {
  return (x + y);
}

...

bits32 a1, b1;
auto c1;
a1 = 5;
b1 = 10;

c1 = add(b1, c1);

bits64 a2, b2;
auto c2;
a2 = 5;
b2 = 10;

c2 = add(b2, c2);
    \end{lstlisting}
    %\begin{minipage}[t]{0.5\linewidth}
    \end{minipage}%
    \begin{minipage}{0.5\linewidth}
    \begin{lstlisting}
// Automatically generated
// monomorphic definitions
_Madd_$f$Lb32...(bits32 x,
                 bits32 y) {
    return (x + y);
}
_Madd_$f$Lb64...(bits64 x,
                 bits64 y) {
    return (x + y);
}

...

bits32 a1, b1;
bits32 c1; // filled-in
a1 = 5;
b1 = 10;

c1 = _Madd_$f$Lb32...(a1, b1);

bits64 a2, b2;
bits64 c2; // filled-in
a2 = 5;
b2 = 10;

c2 = _Madd_$f$Lb64...(a2, b2);
    \end{lstlisting}
    \end{minipage}
    \end{center}
\end{listing}

\subsection{Simplified pointer dereferences}

In \cmm, a reference to memory (used for indirect write on indirect read) has to specify the type of the memory location using the syntax \li{type[address]}. This syntax is close to C array indexing, but it consists of a type parameter and an address parameter (basically, a pointer). Introducing type inference into the language allowed us to make a simplification of this syntax with the type left out. The type is then inferred from the address parameter, or, from the result. Any inferred type of the result or the address parameter then propagates to the other one. We show this in \cref{lst:fieldAccessors}.

\subsection{Translation into LLVM}

The prototype compiler translates the input code into the LLVM assembly language. This requires transforming the code into a simpler, flat representation, with subsequent flow analysis, which then allows the code be transcribed into the SSA representation of LLVM. Performing flow analysis allows for, among other things, identifying unreachable code, which is still type-inferred, but not translated, and identifying uninitialized variables, which is followed by generating an error for each. The example in \cref{lst:translex} shows the translation without unreachable code and showcases an fairly unusual assignment statement, which is a feature of \cmm.

\subsection{Improvements over the Algorithm W}

In our previous work \cite{klepl2020type}, we used an variation of Algorithm W based on THIH \cite{jones1999typing}. We concluded that such approach did not allow for multi-parameter typeclasses (\emph{MPTC}s) and general use of field accessors to \li{struct} fields, we required each field to have a unique name.

In the prototype compiler, we address this and the prototype compiler's type inference interprets both MPTCs and field accessors without requirement on unique names.
\begin{listing}
    \small
    \caption{Example of translation into LLVM}
    \label{lst:translex}
    \begin{center}
    %\begin{minipage}[t]{0.5\linewidth}
    \begin{minipage}{0.5\linewidth}
    \begin{lstlisting}
// Source code
foreign "C" f(bits32 u,
              bits32 w)
              -> bits32{
  auto r;

  if u == w {
    r, u = u, w;
    return (u);
  } else {
    w, r = u, w;
    u = r + w;
    return (u);
  }

  u = u * 2;
  w = u + r;

  return (w);
}
    \end{lstlisting}
    %\begin{minipage}[t]{0.5\linewidth}
    \end{minipage}%
    \begin{minipage}{0.5\linewidth}
    \begin{lstlisting}[language=LLVM]
; Output LLVM assembly
define external ccc  <{i32}>
@f(i32  %u_0, i32  %w_0)    {
LR$procedure0:
  %0 = icmp eq i32 %u_0, %w_0
  br i1 %0, label %F$then_1,
            label %F$else_1
F$then_1:
  %1 = phi i32
        [%u_0, %LR$procedure0]
  %2 = phi i32
        [%w_0, %LR$procedure0]
  %3 = insertvalue <{i32}>
        undef, i32 %2, 0
  ret <{i32}> %3
F$else_1:
  %4 = phi i32
        [%u_0, %LR$procedure0]
  %5 = phi i32
        [%w_0, %LR$procedure0]
  %6 = add   i32 %5, %4
  %7 = insertvalue <{i32}>
        undef, i32 %6, 0
  ret <{i32}> %7
  ; the last three statements
  ; were unreachable
}
    \end{lstlisting}
    \end{minipage}
    \end{center}
\end{listing}

\subsubsection{Multi-parameter typeclasses}
\label{sec:mptcs}

Multi-parameter typeclasses allow us to overload procedures with more than one type argument, this is useful for situations when the procedure is to be implemented differently depending on the parameter types. Perhaps, the procedure can be implemented more efficiently, when the input types are the same. The example in \cref{lst:mptcs} shows exactly that. If both of the types are \li{bits32}, the instance adds the inputs together and returns the result. If the second argument is a pointer, it performs some more complicated operation on the first argument, represented by multiplication.
\begin{listing}
    \small
    \caption{Example of MPTC usage in the extended \cmm (some names truncated for brevity)}
    \label{lst:mptcs}
    \begin{center}
    %\begin{minipage}[t]{0.5\linewidth}
    \begin{minipage}{0.5\linewidth}
    \begin{lstlisting}[basicstyle=\scriptsize\ttfamily]
// Source code
class C a b {
  m (auto (a) x, auto (b) y)
    -> auto (a);
}

instance C bits32 bits32 {
  m (auto u, auto w) {
    w = w + u;
    return (w);
  }
}

instance C bits32 ptr(auto) {
  m (auto u, auto w) {
    u = u * 4;
    return (u);
  }
}

foreign "C"
f(bits32 x, ptr(bits64) y) {
  x = m(x, y);

  return (x);
}

foreign "C"
g(bits32 x) {
  x = m(x, x);

  return (x);
}
    \end{lstlisting}
    %\begin{minipage}[t]{0.5\linewidth}
    \end{minipage}%
    \begin{minipage}{0.5\linewidth}
    \begin{lstlisting}[language=LLVM, basicstyle=\scriptsize\ttfamily]
; Output LLVM assembly
define external ccc  <{i32}>
  @_Mm_$f$Lb32b64...(
    i32  %u_0, i64*  %w_0) {
LR$procedure2:
  %0 = mul   i32 %u_0, 4
  %1 = insertvalue <{i32}>
        undef, i32 %0, 0
  ret <{i32}> %1
}

define external ccc  <{i32}>
  @_Mm_$f$Lb32b32...(
    i32  %u_0, i32  %w_0) {
LR$procedure1:
  %0 = add   i32 %w_0, %u_0
  %1 = insertvalue <{i32}>
        undef, i32 %0, 0
  ret <{i32}> %1
}

define external ccc  <{i32}>
  @f(i32  %x_0, i64*  %y_0) {
LR$procedure3:
  %0 =  call ccc  <{i32}>
          @_Mm_$f$Lb32b64...(
            i32  %x_0,
            i64*  %y_0)
  %1 = extractvalue <{i32}>%0,0
  %2 = insertvalue <{i32}>
        undef, i32 %1, 0
  ret <{i32}> %2
}

define external ccc  <{i32}>
  @g(i32  %x_0) {
LR$procedure4:
  %0 =  call ccc  <{i32}>
          @_Mm_$f$Lb32b32...(
            i32  %x_0,
            i32  %x_0)
  %1 = extractvalue <{i32}>%0,0
  %2 = insertvalue <{i32}>
        undef, i32 %1, 0
  ret <{i32}> %2
}
    \end{lstlisting}
    \end{minipage}
    \end{center}
\end{listing}

\subsubsection{User-specified functional dependencies}
\label{sec:uspecFD}

Sometimes, a procedure can have returns of different types depending on the given arguments. In many instances, these types are closely related to the argument types and they can be derived from type algebra. For such cases, regular (multi-parameter) typeclasses are enough as demonstrated by other examples. But, sometimes the types of the returned values are not algebraicly dependent on the argument types. One example of this are the oftenly mentioned field accessors.

The prototype compiler allows definition of functional dependencies for arbitrary user-defined typeclass. A simple use of this is demonstrated in \cref{lst:fundeps}. The compiler correcly type-infers the result types according to the given functional dependency.

\begin{listing}
    \caption{Example use of functional dependencies (some names truncated for brevity)}
    \label{lst:fundeps}
    \begin{center}
    %\begin{minipage}[t]{0.5\linewidth}
    \begin{minipage}{0.5\linewidth}
    \begin{lstlisting}[basicstyle=\scriptsize\ttfamily]
// Source code

// a type class with
//  a fun. dependency: a -> b
class Id a b | a -> b {
  id (auto (a) x) -> auto (b);
}

instance Id bits32 bits32 {
  id (auto u) {
    return (u);
  }
}

instance Id bits64 bits64 {
  id (auto u) {
    return (u);
  }
}

...

bits32 x;
auto id_x;
bits64 y;
auto id_y;

id_x = id(x);
id_y = id(y);
    \end{lstlisting}
    %\begin{minipage}[t]{0.5\linewidth}
    \end{minipage}%
    \begin{minipage}{0.5\linewidth}
    \begin{lstlisting}[basicstyle=\scriptsize\ttfamily]
// Monomorphized code
//  with interpreted
//  functional dependencies
_Mid_$f$Lb32...(bits32 u) {
  return (u);
}
_Mid_$f$Lb64...(bits64 u) {
  return (u);
}

...

bits32 x;
// type-inferred return
bits32 id_x;

bits64 y;
// type-inferred return
bits64 id_y;

id_x = _Mid_$f$Lb32...(x);
id_y = _Mid_$f$Lb64...(y);
    \end{lstlisting}
    \end{minipage}
    \end{center}
\end{listing}

\subsubsection{Field accessors}

Field accessors allow us to introduce further semantics into structs. This is useful in situations, where each object in a certain model has to contain some information that characterizes the given object. It can be the object's identifier, or age, and perhaps, it can have different types for different objects. For example, one part of the model can use 32-bit wide identifiers and another 64-bit wide. This is where the approach from the previous work did not suffice.

With the prototype implementation based on Deferred Solving with MPTCs and functional dependencies, type-inferring general field accessors is possible as shown in \cref{lst:fieldAccessors}. The compiler type-infers all relevant types, only \li{struct} definitions are typed explicitly.

\begin{listing}
    \caption{Example use of field accessors (some names truncated for brevity)}
    \label{lst:fieldAccessors}
    \begin{center}
    %\begin{minipage}[t]{0.5\linewidth}
    \begin{minipage}{0.45\linewidth}
    \begin{lstlisting}[basicstyle=\footnotesize\ttfamily]
// Source code
struct A {
  id: bits32;
}

struct B {
  id: bits64;
}

update_id(auto s,
          auto x) {
  // indirect write
  // to 'id' field
  [s->id] = x;
  return ();
}

access_id(auto s) {
  auto x;
  // indirect read
  x = [s->id];
  return (x);
}

...

stackdata {
  a: A;
  b: B;
}

auto x;
auto y;

= update_id(a, 5);
= update_id(b, 10);

x = access_id(a);
y = access_id(b);
    \end{lstlisting}
    %\begin{minipage}[t]{0.5\linewidth}
    \end{minipage}%
    \begin{minipage}{0.55\linewidth}
    \begin{lstlisting}[language=LLVM,basicstyle=\scriptsize\ttfamily]
; Output LLVM assembly
define external ccc  <{}>
  @_Mupdate...A({i32}*  %s_0,
                i32  %x_0) {
LR$procedure0:
  %0 = getelementptr {i32}, {i32}* %s_0,
                      i32 0, i32 0
  store   i32 %x_0, i32* %0
  ret <{}> <{  }>
}
; define @_Mupdate...B <- did not fit

define external ccc  <{i32}>
  @_Maccess...A({i32}*  %s_0) {
LR$procedure1:
  %0 = getelementptr {i32}, {i32}* %s_0,
                      i32 0, i32 0
  %1 = load   i32, i32* %0
  %2 = insertvalue <{i32}>
        undef, i32 %1, 0
  ret <{i32}> %2
}
; define @_Maccess...B <- did not fit

...

%0 = alloca {i32}
%1 = alloca {i64}
%2 = call ccc  <{}>
       @_Mupdate...A({i32}* %0, i32 5)
%3 = call ccc  <{}>
       @_Mupdate...B({i64}* %1, i64 10)
%4 = call ccc  <{i32}>
       @_Maccess...A({i32}* %0)
%5 = extractvalue <{i32}> %4, 0
%6 = call ccc  <{i64}>
       @_Maccess...B({i64}* %1)
%7 = extractvalue <{i64}> %6, 0

    \end{lstlisting}
    \end{minipage}
    \end{center}
\end{listing}

\subsection{Typeclass-based resource management}
\label{sec:tcResMan}

Let us say, we want to perform some resource managemet action on every procedure exit. Perhaps, we opened a file and we want to close it. In practice, this is quite common, and it let to, for example, developing RAII in C++, or try/finally pattern in C\#. For this purpose, in the prototype compiler, we can use \li{new}-specialized stack-allocated objects.

If an object carries the \li{new} specialization, the compiler calls \li{drop} procedure on the pointer reference to the object before every exit from the procedure. This \li{drop} call then carries out the resource management action as defined for the object's type. This is demonstrated by \cref{lst:resource}

\begin{listing}
    \small
    \caption{Example use of automatic resource management (some names truncated for brevity)}
    \label{lst:resource}
    \begin{center}
    %\begin{minipage}[t]{0.5\linewidth}
    \begin{minipage}{0.35\linewidth}
    \begin{lstlisting}
// Source code
struct X {
  x: bits32;
}

struct Y {
  x: bits64;
}

class Managed a {
drop(ptr(auto(a)) u)
  -> ;
}

instance Managed X {
  drop(auto u) {
    auto y;
    y = [u->x];
    return ();
  }
}

instance Managed Y {
  drop(auto w) {
    return ();
  }
}

foreign "C"
f(bits32 a) {
  stackdata {
    F: new X;
    G: new Y;
  }

  return ();
}
    \end{lstlisting}
    %\begin{minipage}[t]{0.5\linewidth}
    \end{minipage}%
    \begin{minipage}{0.65\linewidth}
    \begin{lstlisting}[language=LLVM]
; Output LLVM assembly
define external ccc <{}>
@_Mdrop...X({i32}*  %u_0) {
LR$procedure1:
  %0 = getelementptr {i32}, {i32}* %u_0,
                      i32 0, i32 0
  %1 = load   i32, i32* %0
  ret <{}> <{  }>
}


define external ccc <{}>
@_Mdrop...Y({i64}*  %w_0) {
LR$procedure2:
  ret <{}> <{  }>
}


define external ccc <{}>
@f(i32  %a_0) {
LR$procedure3:
  %0 = alloca {i32}
  %1 = alloca {i64}

  ; automatically generated
  ; resource management:

  ; represented by X
  %2 =  call ccc  <{}>
          @_Mdrop...X({i32}*  %0)
  ; represented by Y
  %3 =  call ccc  <{}>
          @_Mdrop...Y({i64}*  %1)
  ret <{}> <{  }>
}
    \end{lstlisting}
    \end{minipage}
    \end{center}
\end{listing}

\subsubsection{Type-inferred resource management}

It may happen the resource management action depends on a type parameter. Perhaps due to architectural specifics of the target platform, or due to specifics of a certain component. For this, we can combine automatic resource management and type inference of the resource objects (with or without typeclass overloading) as demonstrated by \cref{lst:polyResource}

\begin{listing}
    \small
    \caption{Example use of type-inferred resource management (some names truncated for brevity)}
    \label{lst:polyResource}
    \begin{center}
    %\begin{minipage}[t]{0.5\linewidth}
    \begin{minipage}{0.35\linewidth}
    \begin{lstlisting}
// Source code
struct X {
  x: bits32;
}

struct Y {
  x: bits64;
}

drop(auto u) {
  // ...
  return ();
}

foreign "C"
f(ptr(X) s,
  ptr(Y) t) {
  stackdata {
    // type-inferred
    //   resource
    //   objects
    x: new auto;
    y: new auto;
  }

  [x] = [s->x];
  [y] = [t->x];

  return ();
}
    \end{lstlisting}
    %\begin{minipage}[t]{0.5\linewidth}
    \end{minipage}%
    \begin{minipage}{0.65\linewidth}
    \begin{lstlisting}[language=LLVM]
; Output LLVM assembly
define external ccc <{}>
@_Mdrop...(i32*  %u_0) {
LR$procedure0:
  ; ...
  ret <{}> <{  }>
}


define external ccc <{}>
@_Mdrop...(i64*  %u_0) {
LR$procedure0:
  ; ...
  ret <{}> <{  }>
}


define external ccc <{}>
@f({i32}*  %s_0, {i64}*  %t_0) {
LR$procedure1:
  %0 = alloca i32
  %1 = alloca i64
  %2 = getelementptr {i32}, {i32}* %s_0,
                      i32 0, i32 0
  %3 = load   i32, i32* %2
  store   i32 %3, i32* %0
  %4 = getelementptr {i64}, {i64}* %t_0,
                      i32 0, i32 0
  %5 = load   i64, i64* %4
  store   i64 %5, i64* %1
  %6 =  call ccc  <{}>
          @_Mdrop...(i32*  %0)
  %7 =  call ccc  <{}>
          @_Mdrop...(i64*  %1)
  ret <{}> <{  }>
}
    \end{lstlisting}
    \end{minipage}
    \end{center}
\end{listing}

\subsection{Automatic resource management with \texttt{dropped} statements}

Sometimes, when we use automatic resource management, there exists a special case, in which we do not want to perform the resource management action. This may apply to cases where the resource is a file handle of a file we failed to open. In the prototype compiler, we can use \li{dropped} statement with a named reference to a resource object to specify that the corresponding resource management action is not to be performed automatically (but may be still manually). We demonstrate its use in \cref{lst:manual}.


\begin{listing}
    \small
    \caption{Example use of automatic resource management with \li{dropped} statements (some names truncated for brevity)}
    \label{lst:manual}
    \begin{center}
    %\begin{minipage}[t]{0.5\linewidth}
    \begin{minipage}{0.35\linewidth}
    \begin{lstlisting}
// Source code
drop(auto a) {
  return ();
}

foreign "C"
f(bits32 a) {
  stackdata {
    x:y: new bits32;
  }

  if a == 0 {
    dropped x;
    // manual drop
    // (then_1)
    = drop(x);
    return ();
  }

  if a == 5 {
    if a == 10 {
  // automatic drop
  // (then_3)
      return ();
    } else {
      dropped y;
      // no drop
      return ();
    }
  }

  // automatic drop
  // (fi_2)
  return ();
}
    \end{lstlisting}
    %\begin{minipage}[t]{0.5\linewidth}
    \end{minipage}%
    \begin{minipage}{0.65\linewidth}
    \begin{lstlisting}[language=LLVM]
; Output LLVM assembly
define external ccc <{}>
@_Mdrop...(i32*  %a_0) {
LR$procedure0:
  ret <{}> <{  }>
}


define external ccc <{}>
@f(i32  %a_0) {
LR$procedure1:
  %0 = alloca i32
  %1 = icmp eq i32 %a_0, 0
  br i1 %1, label %F$then_1,
            label %F$fi_1
F$then_1:
  %2 =  call ccc  <{}>
          @_Mdrop...(i32*  %0)
  ret <{}> <{  }>
F$fi_1:
  %3 = phi i32 [%a_0, %LR$procedure1]
  %4 = icmp eq i32 %3, 5
  br i1 %4, label %F$then_2,
            label %F$fi_2
F$then_2:
  %5 = phi i32 [%3, %F$fi_1]
  %6 = icmp eq i32 %5, 10
  br i1 %6, label %F$then_3,
            label %F$else_3
F$fi_2:
  %7 =  call ccc  <{}>
          @_Mdrop...(i32*  %0)
  ret <{}> <{  }>
F$then_3:
  %8 =  call ccc  <{}>
          @_Mdrop...(i32*  %0)
  ret <{}> <{  }>
F$else_3:
  ret <{}> <{  }>
}
    \end{lstlisting}
    \end{minipage}
    \end{center}
\end{listing}

\section{Applicability to other languages}

\subsection{Subtyping in systems programming}

The subset of subtyping modelled by our proposed type system can be used to check and infer various secondary properties of objects. Throughout the thesis, we demonstrated its use on constnesses and data kinds of the \cmm language. Connecting the analysis of these secondary properties to the type system allows for tying the constraints on these properties only to a particular typing (as defined in \cref{sec:typesystem}), for example, by defining each typeclass instance with a different set of constraints.

\subsection{C subtyping}

In the previous work \cite{klepl2020type}, we identified a weakness of the then-used type system with regards to the \li{const} qualifier from the C language \cite{cstandard2018}. Let us examine whether its semantics fit the type system we introduced in \cref{sec:typesystem}.

The semantics of the C \li{const} qualifier (as specified by \cite{cstandard2018}), quite differ from the constness introduced in \cref{const_kind_def}, despite the similarity of names. For our purposes, we simplify the semantics into two cases:

\begin{enumerate}
    \item If the \li{const} qualifier is applied to a variable (for example, \li{char *const p}), the type of the variable is the same as if it was not \li{const}-qualified. Any attempt to assign to this variable is an undefined behavior. \label{constvar}

    \item If the \li{const} qualifier is applied to a type pointed to by a certain pointer (for example, \li{char const *p}), any assignments to the pointer from pointers pointing to a non-\li{const}-qualified version of the same type are valid and the two pointers have the same value afterwards. The reverse assignment would be, again, an undefined behavior. \label{pointerSubtype}

    In our formalism, this means that \li{T *} is less general than \li{T const *}.
\end{enumerate}

We can model the semantics of the case number \ref{pointerSubtype} (case number \ref{constvar} is identified syntactically) via introducing two new subtype dimensions, one which is only relevant to pointer types (in simplification; we could then generalize this to array types and other aggregate types), let us call this dimension the \emph{C-const dimension}. And one other dimension to represent all subtype relations, let us call this dimension the \emph{assigment dimension}.

\begin{defn}[C-const dimension]
    The C-const subtype dimension $C$ is defined by two constants $\mtt{NonConst}, \mtt{Const}$ and an ordering $\top = \mtt{Const} >_C \mtt{NonConst} = \bot$.

    This dimension is only relevant to pointer types, other types being either unconstrained or constrained to specifically one of the constants (which is then sometimes called canonical -- see \cite{tiuryn1999subtyping}).
\end{defn}

\begin{defn}[Assignment dimension]
    The assignment subtype dimension $A$ is defined by one constant $\bot_A$.
\end{defn}

In the following, we use $a \sim_\chi b$ as a shortcut for $a \leq_\chi b \land b \leq_\chi a$, where $\chi$ is a dimension of the type system.

\begin{itemize}
    \item Each type definition (or primitive type: for example, \li{int}) $t$ is constrained: $t \sim_C \mtt{NonConst}$ (on top of any other constraints).

    \item Any declaration \li{T const *} constrains the type $t$ of the declared variable: $(\forall \mtt{dimension}\ s, s \neq C. t \sim_s T\ast), t \sim_C \mtt{Const}$ (this models that applying \li{const} is idempotent).

    \item For (generalized) assignments with operand types $t_l, t_r$, we put $t_l \geq_s t_r$ for each dimension $s$ and $t_l \sim_T t_r$. (we say ``generalized assignments'' as we consider even actual arguments to procedures and return operands, etc.).
\end{itemize}

To sufficiently constrain the space so the modelled semantics are the same as the semantics of the C \li{const} qualifier, after deciding all typings and before deciding any subtype dimensions, we introduce new constraints.

These constraints are derived from the assignment dimension $A$ for each two types such that $t \leq_A t'$. For example, if $t \equiv T\ast\ast$ and $t' \equiv T'\ast\ast$, we put $T \sim T'$ (as this is still part of the semantic case number \ref{pointerSubtype}). We will not go into further details as this is outside the scope of this work. After this step, the type inference continues as usual.

We leave this for future work together with considering the possibility of generating all subtype constraints after deciding the typings for the given subprogram component.

\subsection{Automatic resource management in practical systems programming}
\label{sec:autoC}

We show how to use the results of this thesis to simplify the resource management in system programming languages.

The kernel coding style guide proposes the following convention:

\begin{quotation}[\xxx{add pointer to kernel group}]
  The goto statement comes in handy when a procedure exits from multiple locations and some common work such as cleanup has to be done. If there is no cleanup needed then just return directly.
\end{quotation}

We propose the automatic resource management features introduced by the extended \cmm as an alternative to this convention. The cleanup section is replaced by a definition of an resource-representing object and accompanying \li{drop} procedure.

In order to fully serve as a practical alternative to the \li{goto cleanup} pattern in a general case, we also propose the use of a special statement \li{dropped}. Our introduced \li{dropped} statement takes a named alias to the object representing the resource and removes the implicit calls on all procedure exits following this statement, still allowing for explicit calls whenever the user specifies.

One possible improvement over the \li{goto cleanup} pattern is that this can easily scale with more actions tied to multiple resource-representing objects, each such action then possibly being prevented by a separate \li{dropped} statement on some subset of control flow paths.

\section{Comparison to other languages offering automatic resource management}
\label{sec:OtherRAII}

There are multiple constructs that implement the automatic resource management mechanism. We will look into two: C++ RAII and C\# \li{try\{...\} finally\{...\}} (for simplicity, we will call them by the respective language's names).

Neither of approaches of C++ and C\# allow for explicit use-specified special cases to the resource management mechanism. If the given language has no predefined optimization for the specific situation (C++, notably, defines ``trivially destructible'' objects, which do not issue the resource management action), the resource management action has to happen regardless of circumstances (as long as the program follows the language specification).

One other differences is that both approaches define a symmetric preparation action for the resource management action we consider. C++ calls the resource management action \emph{destructor} and the symmetric action \emph{constructor} (constructors can generally have more possible cases). C\# then puts the preparation action inside the \li{try\{...\}} block of the pattern.

We do not consider this symmetric preparation action as it is usually specific to the given instance (for a file handle, specified the given file name) and the functions have one entry, giving only constant benefit for each procedure, if any, to considering the symmetric action.


\chapwithtoc{Conclusion}

In this thesis, we have described and implemented a rather complex extension of the deferred type inference algorithm, and applied it for inference of types of an extended variant of the \cmm language that features types and a simple but powerful type-driven form of resource management. The output of the thesis, a prototype compiler, translates the typed \cmm language into lower-level, simpler LLVM assembly code, which can be further used to produce executables and libraries.

In particular, we have successfully achieved the following main goals:

\begin{description}
    \small
    \item[Extend \cmm with parametric polymorphism with typeclasses] We achieved this goal successfully and we demonstrate parametric polymorphism with typeclasses in our results, see \cref{sec:reduction}. We also extended the typeclasses to multi-parameter typeclasses with user-specified functional dependencies, solving one of the issues identified in the previous work \cite{klepl2020type}. This allows for more expressive code and stronger type constraints than regular typeclasses, see \cref{sec:mptcs,sec:uspecFD}.

    \item[Extend \cmm with record types]  We extended the language with \li{struct} types (\cref{sec:structExt}) inspired by C, but allowing for multiple aliases to fields. The field accessors are defined as singleton methods of classes parametrized by corresponding structs and their respective fields and with functional dependencies \li{struct -> field}, which solves the field-accessor-related issue identified in the previous work.

    \item[Extend \cmm with automatic resource management features] We extended the language with \li{new} specifier (\cref{RAII}), which signifies that an object is to automatically undergo a predefined resource management action on each procedure exit. We also successfully experimented with introducing user-specified special cases to this mechanism and with various combinations with the type inference mechanism. We showed practical uses-cases for each of these further extensions, see \cref{sec:tcResMan}.

    \item[The solution should simplify implementation of low-level software] We showed that the solution can be used to reduce the written code (\cref{sec:reduction}) and we also discussed the effects on performance and showed the generic code does not constitute any runtime type checks or indirections.

    \item[The extended \cmm  code compiles into LLVM] In \cref{sec:protoComp}, we showed many examples of the input code and the corresponding translation to LLVM. In \cref{sec:flatBlock,sec:typeInf,sec:mono,translation}, we explained the steps necessary for such translation and discussed some optimizations that use the same primitives as the translation and which can be performed simultaneously with the semantic processing of the code for translation. The compiler also contains options \li{-f} and \li{-n} for running the compiler pipeline up to the flattening phase and just before the transcription to LLVM, respectively \cite{klepl2022compiler}. These options help connect the input code to its LLVM representation.

    \item[The available features will be evaluated by demonstrational programs] In our results (\cref{chap4}), we include an example for every claim we make about the extended \cmm in results or here. Furthermore, there are even more examples, that can be automatically executed by a script, in the supplementary repository \cite{klepl2022compiler}.

    \item[The available features will be compared with similarly-focused languages] We compared the features of the extended \cmm to C, the ``de facto'' standard system programming language and we discussed possible benefits of our approach to some common design patterns of C (\cref{sec:autoC}).

    We also compared the features of the extended \cmm to other languages offering support for generics, namely Java and Python. And we showed that our solution, contrasting these, does not require any runtime type information (\cref{sec:reduction}).

    We then compared the features of the extended \cmm to other languages offering automatic resource management, namely C++ and C\#. And discussed the differences in design decisions (\cref{sec:OtherRAII}).

    Less-significantly, we compared the type system introduced into \cmm to the one we introduced into C in the previous work and we concluded that it solves the unreached goals of that work.

    \item[User front-end for customizable subtyping systems] While, in \cref{sec:typesystem}, we have shown that the subtyping system is capable of inferring user-specified semantics of subtypes -- for example, giving a possibility for the user to simulate various levels of constantness known from C++ (\li{constexpr}), we did not develop a practical front-end that would allow programmers to interact with this feature directly.
\end{description}

% \section*{The overall quality of the solution}

% In the prototype compiler, we focused on the assignment goals and on the goals the previous work did not manage to achieved. We also addressed concerns with stability and usability of the compiler included in the previous work. These concerns were mostly about compiler errors that apparently did not reflect the actual errors made by the user and about platform-specificity of the solution. The functionality of the compiler of the previous work was tied to some platform specifics which were not sufficiently reflected by the provided documentation. In the prototype compiler for this thesis, we address both these issues. The functionality of the compiler does not put any assumptions on the actual platform other than the code dependencies stated in included documentation. The concern with errors is addressed by introducing \li{ASTErrors} extensively used throughout the implementation, that on top of the specific also print out the assumed relevant location in the code that issues the error.

% \section*{Practical uses of the solution}

% We discussed practical use-cases of our solution in systems programming, specifically, in Linux kernel programming. We also discussed various general use-cases for polymorphism with typeclasses for reducing human-written by introducing polymorphic definitions and for automatic resource management. We supported these discussions by examples compiled using the prototype compiler and compared them to similar solutions of popular languages.

\section*{Future work}

Here we include the possible extensions to our solution that were outside of time constraints of our work.

\begin{description}
    \item[Existential types] We based our solution for type inference on deferred solving algorithm capable of solving existential types \cite{vytiniotis2011outsidein}. One natural course of future development could be trying to combine those with our extensions to the language. Usefulness of existential types in imperative languages has been demonstrated, for example, by \citet{grossman2002existential}.

    \item[Stronger subtyping system] We discussed some insufficiencies of the type system we developed. We leave for a future work a revised approach, which constitutes two constraint generation phases as described in \cref{sec:weakness}. This revised approach can generate stronger subtype constraints.

    At this point, it is unclear whether the complications induced by better subtype support, mainly the complexity of inference, would be sufficiently counter-weighted by the newly gained features. Further research is needed to show if there is any significant benefit in letting the programmers to interact with the subtype system -- e.g., defining their own semantics and interpretation for constness-like annotations.
\end{description}


\ifEN
\chapwithtoc{Bibliography}
\else
\chapwithtoc{Seznam použité literatury}
\fi

\printbibliography[heading=none]


\appendix
\include{howto}

% if your attachments are complicated, describe them in a separate appendix
%\include{attachments}

\openright
\end{document}
